%% process.tex
%%
%% Copyright (C) 2000-2018 by Simone Piccardi.  Permission is granted to
%% copy, distribute and/or modify this document under the terms of the GNU Free
%% Documentation License, Version 1.1 or any later version published by the
%% Free Software Foundation; with the Invariant Sections being "Un preambolo",
%% with no Front-Cover Texts, and with no Back-Cover Texts.  A copy of the
%% license is included in the section entitled "GNU Free Documentation
%% License".
%%

\chapter{L'interfaccia base con i processi}
\label{cha:process_interface}

Come accennato nell'introduzione il \textsl{processo} è l'unità di base con
cui un sistema unix-like alloca ed utilizza le risorse.  Questo capitolo
tratterà l'interfaccia base fra il sistema e i processi, come vengono passati
gli argomenti, come viene gestita e allocata la memoria, come un processo può
richiedere servizi al sistema e cosa deve fare quando ha finito la sua
esecuzione. Nella sezione finale accenneremo ad alcune problematiche generiche
di programmazione.

In genere un programma viene eseguito quando un processo lo fa partire
eseguendo una funzione della famiglia \func{exec}; torneremo su questo e sulla
creazione e gestione dei processi nel prossimo capitolo. In questo
affronteremo l'avvio e il funzionamento di un singolo processo partendo dal
punto di vista del programma che viene messo in esecuzione.


\section{Esecuzione e conclusione di un programma}

Uno dei concetti base di Unix è che un processo esegue sempre uno ed un solo
programma: si possono avere più processi che eseguono lo stesso programma ma
ciascun processo vedrà la sua copia del codice (in realtà il kernel fa sì che
tutte le parti uguali siano condivise), avrà un suo spazio di indirizzi,
variabili proprie e sarà eseguito in maniera completamente indipendente da
tutti gli altri. Questo non è del tutto vero nel caso di un programma
\textit{multi-thread}, ma la gestione dei \textit{thread} in Linux sarà
trattata a parte\unavref{in cap.~\ref{cha:threads}}.


\subsection{L'avvio e l'esecuzione di un programma}
\label{sec:proc_main}

\itindbeg{link-loader}
\itindbeg{shared~objects}
Quando un programma viene messo in esecuzione, cosa che può essere fatta solo
con una funzione della famiglia \func{exec} (vedi sez.~\ref{sec:proc_exec}),
il kernel esegue un opportuno codice di avvio, il cosiddetto
\textit{link-loader}, costituito dal programma \cmd{ld-linux.so}. Questo
programma è una parte fondamentale del sistema il cui compito è quello della
gestione delle cosiddette \textsl{librerie condivise}, quelle che nel mondo
Windows sono chiamate DLL (\textit{Dinamic Link Library}), e che invece in un
sistema unix-like vengono chiamate \textit{shared objects}.

Infatti, a meno di non aver specificato il flag \texttt{-static} durante la
compilazione, tutti i programmi in Linux sono compilati facendo riferimento a
librerie condivise, in modo da evitare di duplicare lo stesso codice nei
relativi eseguibili e consentire un uso più efficiente della memoria, dato che
il codice di uno \textit{shared objects} viene caricato in memoria dal kernel
una sola volta per tutti i programmi che lo usano.
\itindend{shared~objects}

Questo significa però che normalmente il codice di un programma è incompleto,
contenendo solo i riferimenti alle funzioni di libreria che vuole utilizzare e
non il relativo codice. Per questo motivo all'avvio del programma è necessario
l'intervento del \textit{link-loader} il cui compito è caricare in memoria le
librerie condivise eventualmente assenti, ed effettuare poi il collegamento
dinamico del codice del programma alle funzioni di libreria da esso utilizzate
prima di metterlo in esecuzione.

Il funzionamento di \cmd{ld-linux.so} è controllato da alcune variabili di
ambiente e dal contenuto del file \conffile{/etc/ld.so.conf} che consentono di
elencare le directory un cui cercare le librerie e determinare quali verranno
utilizzate.  In particolare con la variabile di ambiente
\envvar{LD\_LIBRARY\_PATH} si possono indicare ulteriori directory rispetto a
quelle di sistema in cui inserire versioni personali delle librerie che hanno
la precedenza su quelle di sistema, mentre con la variabile di ambiente
\envvar{LD\_PRELOAD} si può passare direttamente una lista di file di librerie
condivise da usare al posto di quelli di sistema. In questo modo è possibile
effettuare lo sviluppo o il test di nuove librerie senza dover sostituire
quelle di sistema. Ulteriori dettagli sono riportati nella pagina di manuale
di \cmd{ld.so} e per un approfondimento dell'argomento si può consultare
sez.~3.1.2 di \cite{AGL}.

Una volta completate le operazioni di inizializzazione di \cmd{ld-linux.so}, il
sistema fa partire qualunque programma chiamando la funzione \code{main}. Sta
al programmatore chiamare così la funzione principale del programma da cui si
suppone che inizi l'esecuzione. In ogni caso senza questa funzione lo stesso
\textit{link-loader} darebbe luogo ad errori.  Lo standard ISO C specifica che
la funzione \code{main} può non avere argomenti o prendere due argomenti che
rappresentano gli argomenti passati da linea di comando (su cui torneremo in
sez.~\ref{sec:proc_par_format}), in sostanza un prototipo che va sempre bene è
il seguente:
\includecodesnip{listati/main_def.c}

\itindend{link-loader}

In realtà nei sistemi Unix esiste un altro modo per definire la funzione
\code{main}, che prevede la presenza di un terzo argomento, \code{char
  *envp[]}, che fornisce l'\textsl{ambiente} del programma; questa forma però
non è prevista dallo standard POSIX.1 per cui se si vogliono scrivere
programmi portabili è meglio evitarla. Per accedere all'ambiente, come vedremo
in sez.~\ref{sec:proc_environ} si usa in genere una variabile globale che
viene sempre definita automaticamente.

Ogni programma viene fatto partire mettendo in esecuzione il codice contenuto
nella funzione \code{main}, ogni altra funzione usata dal programma, che sia
ottenuta da una libreria condivisa, o che sia direttamente definita nel
codice, dovrà essere invocata a partire dal codice di \code{main}. Nel caso di
funzioni definite nel programma occorre tenere conto che, nel momento stesso
in cui si usano le librerie di sistema (vale a dire la \acr{glibc}) alcuni
nomi sono riservati e non possono essere utilizzati. 

In particolare sono riservati a priori e non possono essere mai ridefiniti in
nessun caso i nomi di tutte le funzioni, le variabili, le macro di
preprocessore, ed i tipi di dati previsti dallo standard ISO C. Lo stesso
varrà per tutti i nomi definiti negli \textit{header file} che si sono
esplicitamente inclusi nel programma (vedi sez.~\ref{sec:proc_syscall}), ma
anche se è possibile riutilizzare nomi definiti in altri \textit{header file}
la pratica è da evitare nella maniera più assoluta per non generare ambiguità.

Oltre ai nomi delle funzioni di libreria sono poi riservati in maniera
generica tutti i nomi di variabili o funzioni globali che iniziano con il
carattere di sottolineato (``\texttt{\_}''), e qualunque nome che inizi con il
doppio sottolineato (``\texttt{\_\_}'') o con il sottolineato seguito da
lettera maiuscola. Questi identificativi infatti sono utilizzati per i nomi
usati internamente in forma privata dalle librerie, ed evitandone l'uso si
elimina il rischio di conflitti.

Infine esiste una serie di classi di nomi che sono riservati per un loro
eventuale uso futuro da parte degli standard ISO C e POSIX.1, questi in teoria
possono essere usati senza problemi oggi, ma potrebbero dare un conflitto con
una revisione futura di questi standard, per cui è comunque opportuno
evitarli, in particolare questi sono:
\begin{itemize*}
\item i nomi che iniziano per ``\texttt{E}'' costituiti da lettere maiuscole e
  numeri, che potrebbero essere utilizzati per nuovi codici di errore (vedi
  sez.~\ref{sec:sys_errors}),
\item i nomi che iniziano con ``\texttt{is}'' o ``\texttt{to}'' e costituiti
  da lettere minuscole che potrebbero essere utilizzati da nuove funzioni per
  il controllo e la conversione del tipo di caratteri,
\item i nomi che iniziano con ``\texttt{LC\_}'' e costituiti
  da lettere maiuscole che possono essere usato per macro attinenti la
  localizzazione,% mettere in seguito (vedi sez.~\ref{sec:proc_localization}),
\item nomi che iniziano con ``\texttt{SIG}'' o ``\texttt{SIG\_}'' e costituiti
  da lettere maiuscole che potrebbero essere usati per nuovi nomi di segnale
  (vedi sez.~\ref{sec:sig_classification}),
\item nomi che iniziano con ``\texttt{str}'', ``\texttt{mem}'', o
  ``\texttt{wcs}'' e costituiti da lettere minuscole che possono essere
  utilizzati per funzioni attinenti la manipolazione delle stringhe e delle
  aree di memoria,
\item nomi che terminano in ``\texttt{\_t}'' che potrebbero essere utilizzati
  per la definizione di nuovi tipi di dati di sistema oltre quelli di
  tab.~\ref{tab:intro_primitive_types}).
\end{itemize*}


\subsection{Chiamate a funzioni e \textit{system call}}
\label{sec:proc_syscall}

Come accennato in sez.~\ref{sec:intro_syscall} un programma può utilizzare le
risorse che il sistema gli mette a disposizione attraverso l'uso delle
opportune \textit{system call}. Abbiamo inoltre appena visto come all'avvio un
programma venga messo in grado di chiamare le funzioni fornite da eventuali
librerie condivise da esso utilizzate. 

Vedremo nel resto della guida quali sono le risorse del sistema accessibili
attraverso le \textit{system call} e tratteremo buona parte delle funzioni
messe a disposizione dalla libreria standard del C, in questa sezione però si
forniranno alcune indicazioni generali sul come fare perché un programma possa
utilizzare queste funzioni.

\itindbeg{header~file}

In sez.~\ref{sec:intro_standard} abbiamo accennato come le funzioni definite
nei vari standard siano definite in una serie di \textit{header file} (in
italiano \textsl{file di intestazione}).  Vengono chiamati in questo modo quei
file, forniti insieme al codice delle librerie, che contengono le
dichiarazioni delle variabili, dei tipi di dati, delle macro di preprocessore
e soprattutto delle funzioni che fanno parte di una libreria.

Questi file sono necessari al compilatore del linguaggio C per ottenere i
riferimenti ai nomi delle funzioni (e alle altre risorse) definite in una
libreria, per questo quando si vogliono usare le funzioni di una libreria
occorre includere nel proprio codice gli \textit{header file} che le
definiscono con la direttiva \code{\#include}. Dato che le funzioni devono
essere definite prima di poterle usare in genere gli \textit{header file}
vengono inclusi all'inizio del programma. Se inoltre si vogliono utilizzare le
macro di controllo delle funzionalità fornite dai vari standard illustrate in
sez.~\ref{sec:intro_gcc_glibc_std} queste, come accennato, dovranno a loro
volta essere definite prima delle varie inclusioni.

Ogni libreria fornisce i propri file di intestazione per i quali si deve
consultare la documentazione, ma in tab.~\ref{tab:intro_posix_header} si sono
riportati i principali \textit{header file} definiti nella libreria standard
del C (nel caso la \acr{glibc}) che contengono le varie funzioni previste
negli standard POSIX ed ANSI C, e che prevedono la definizione sia delle
funzioni di utilità generica che delle interfacce alle \textit{system call}. In
seguito per ciascuna funzione o \textit{system call} che tratteremo
indicheremo anche quali sono gli \textit{header file} contenenti le necessarie
definizioni.

\begin{table}[htb]
  \footnotesize
  \centering
  \begin{tabular}[c]{|l|c|c|l|}
    \hline
    \multirow{2}{*}{\textbf{Header}}&
    \multicolumn{2}{|c|}{\textbf{Standard}}&
    \multirow{2}{*}{\textbf{Contenuto}} \\
    \cline{2-3}
    & ANSI C& POSIX& \\
    \hline
    \hline
    \headfiled{assert.h}&$\bullet$&    --   & Verifica le asserzioni fatte in un
                                              programma.\\ 
    \headfiled{ctype.h} &$\bullet$&    --   & Tipi standard.\\
    \headfiled{dirent.h}&   --    &$\bullet$& Manipolazione delle directory.\\
    \headfiled{errno.h} &   --    &$\bullet$& Errori di sistema.\\
    \headfiled{fcntl.h} &   --    &$\bullet$& Controllo sulle opzioni dei
                                              file.\\ 
    \headfiled{limits.h}&   --    &$\bullet$& Limiti e parametri del sistema.\\
    \headfiled{malloc.h}&$\bullet$&    --   & Allocazione della memoria.\\
    \headfiled{setjmp.h}&$\bullet$&    --   & Salti non locali.\\
    \headfiled{signal.h}&   --    &$\bullet$& Gestione dei segnali.\\
    \headfiled{stdarg.h}&$\bullet$&    --   & Gestione di funzioni a argomenti
                                             variabili.\\ 
    \headfiled{stdio.h} &$\bullet$&    --   & I/O bufferizzato in standard ANSI
                                              C.\\ 
    \headfiled{stdlib.h}&$\bullet$&    --   & Definizioni della libreria
                                              standard.\\ 
    \headfiled{string.h}&$\bullet$&    --   & Manipolazione delle stringhe.\\
    \headfiled{time.h}  &   --    &$\bullet$& Gestione dei tempi.\\
    \headfiled{times.h} &$\bullet$&    --   & Gestione dei tempi.\\
    \headfiled{unistd.h}&   --    &$\bullet$& Unix standard library.\\
    \headfiled{utmp.h}  &   --    &$\bullet$& Registro connessioni utenti.\\
    \hline
  \end{tabular}
  \caption{Elenco dei principali \textit{header file} definiti dagli standard
    POSIX e ANSI C.}
  \label{tab:intro_posix_header}
\end{table}

Un esempio di inclusione di questi file, preso da uno dei programmi di
esempio, è il seguente, e si noti come gli \textit{header file} possano essere
referenziati con il nome fra parentesi angolari, nel qual caso si indica l'uso
di quelli installati con il sistema,\footnote{in un sistema GNU/Linux che
  segue le specifiche del \textit{Filesystem Hierarchy Standard} (per maggiori
  informazioni si consulti sez.~1.2.3 di \cite{AGL}) si trovano sotto
  \texttt{/usr/include}.}  o fra virgolette, nel qual caso si fa riferimento
ad una versione locale, da indicare con un \textit{pathname} relativo:
\includecodesnip{listati/main_include.c}

Si tenga presente che oltre ai nomi riservati a livello generale di cui si è
parlato in sez.~\ref{sec:proc_main}, alcuni di questi \textit{header file}
riservano degli ulteriori identificativi, il cui uso sarà da evitare, ad
esempio si avrà che:
\begin{itemize*}
\item in \headfile{dirent.h} vengono riservati i nomi che iniziano con
  ``\texttt{d\_}'' e costituiti da lettere minuscole,
\item in \headfile{fcntl.h} vengono riservati i nomi che iniziano con
  ``\texttt{l\_}'', ``\texttt{F\_}'',``\texttt{O\_}'' e ``\texttt{S\_}'',
\item in \headfile{limits.h} vengono riservati i nomi che finiscono in
  ``\texttt{\_MAX}'',
\item in \headfile{signal.h} vengono riservati i nomi che iniziano con
  ``\texttt{sa\_}'' e ``\texttt{SA\_}'',
\item in \headfile{sys/stat.h} vengono riservati i nomi che iniziano con
  ``\texttt{st\_}'' e ``\texttt{S\_}'',
\item in \headfile{sys/times.h} vengono riservati i nomi che iniziano con
  ``\texttt{tms\_}'',
\item in \headfile{termios.h} vengono riservati i nomi che iniziano con
  ``\texttt{c\_}'', ``\texttt{V}'', ``\texttt{I}'', ``\texttt{O}'' e
  ``\texttt{TC}'' e con ``\texttt{B}'' seguito da un numero,
\item in \headfile{grp.h} vengono riservati i nomi che iniziano con
  ``\texttt{gr\_}'',
\item in \headfile{pwd.h} vengono riservati i nomi che iniziano con
  ``\texttt{pw\_}'',
\end{itemize*}

\itindend{header~file}

Una volta inclusi gli \textit{header file} necessari un programma potrà
richiamare le funzioni di libreria direttamente nel proprio codice ed accedere
ai servizi del kernel; come accennato infatti normalmente ogni \textit{system
  call} è associata ad una omonima funzione di libreria, che è quella che si
usa normalmente per invocarla.

Occorre però tenere presente che anche se dal punto di vista della scrittura
del codice la chiamata di una \textit{system call} non è diversa da quella di
una qualunque funzione ordinaria, la situazione è totalmente diversa
nell'esecuzione del programma. Una funzione ordinaria infatti viene eseguita,
esattamente come il codice che si è scritto nel corpo del programma, in
\textit{user space}. Quando invece si esegue una \textit{system call}
l'esecuzione ordinaria del programma viene interrotta con quello che viene
usualmente chiamato un \itindex{context~switch} \textit{context
  switch};\footnote{in realtà si parla più comunemente di \textit{context
    switch} quando l'esecuzione di un processo viene interrotta dal kernel
  (tramite lo \textit{scheduler}) per metterne in esecuzione un altro, ma il
  concetto generale resta lo stesso: l'esecuzione del proprio codice in
  \textit{user space} viene interrotta e lo stato del processo deve essere
  salvato per poterne riprendere l'esecuzione in un secondo tempo.}  il
contesto di esecuzione del processo viene salvato in modo da poterne
riprendere in seguito l'esecuzione ed i dati forniti (come argomenti della
chiamata) vengono trasferiti al kernel che esegue il codice della
\textit{system call} (che è codice del kernel) in \textit{kernel space}; al
completamento della \textit{system call} i dati salvati nel \textit{context
  switch} saranno usati per riprendere l'esecuzione ordinaria del programma.

Dato che il passaggio dei dati ed il salvataggio del contesto di esecuzione
sono operazioni critiche per le prestazioni del sistema, per rendere il più
veloce possibile questa operazione sono state sviluppate una serie di
ottimizzazioni che richiedono alcune preparazioni abbastanza complesse dei
dati, che in genere dipendono dall'architettura del processore e sono scritte
direttamente in \textit{assembler}.


%
% TODO:trattare qui, quando sarà il momento vsyscall e vDSO, vedi:
% http://davisdoesdownunder.blogspot.com/2011/02/linux-syscall-vsyscall-and-vdso-oh-my.html 
% http://www.win.tue.nl/~aeb/linux/lk/lk-4.html
%
% Altro materiale al riguardo http://lwn.net/Articles/615809/
% http://man7.org/linux/man-pages/man7/vdso.7.html 

Inoltre alcune \textit{system call} sono state modificate nel corso degli anni
con lo sviluppo del kernel per aggiungere ad esempio funzionalità in forma di
nuovi argomenti, o per consolidare diverse varianti in una interfaccia
generica.  Per questo motivo dovendo utilizzare una \textit{system call} è
sempre preferibile usare l'interfaccia fornita dalla \textsl{glibc}, che si
cura di mantenere una uniformità chiamando le versioni più aggiornate.

Ci sono alcuni casi però in cui può essere necessario evitare questa
associazione, e lavorare a basso livello con una specifica versione, oppure si
può voler utilizzare una \textit{system call} che non è stata ancora associata
ad una funzione di libreria.  In tal caso, per evitare di dover effettuare
esplicitamente le operazioni di preparazione citate, all'interno della
\textsl{glibc} è fornita una specifica funzione,
\funcd{syscall},\footnote{fino a prima del kernel 2.6.18 per l'esecuzione
  diretta delle \textit{system call} erano disponibili anche una serie di
  macro \texttt{\_syscall\textsl{N}} (con $N$ pari al numero di argomenti
  della \textit{system call}); queste sono deprecate e pertanto non ne
  parleremo ulteriormente.} che consente eseguire direttamente una
\textit{system call}; il suo prototipo, accessibile se si è definita la macro
\macro{\_GNU\_SOURCE}, è:

\begin{funcproto}{
  \fhead{unistd.h}
  \fhead{sys/syscall.h}
  \fdecl{long syscall(int number, ...)}
  \fdesc{Esegue la \textit{system call} indicata da \param{number}.}
}
{La funzione ritorna un intero dipendente dalla \textit{system call} invocata,
 in generale $0$ indica il successo ed un valore negativo un errore.}
\end{funcproto}

La funzione richiede come primo argomento il numero della \textit{system call}
da invocare, seguita dagli argomenti da passare alla stessa, che ovviamente
dipendono da quest'ultima, e restituisce il codice di ritorno della
\textit{system call} invocata. In generale un valore nullo indica il successo
ed un valore negativo è un codice di errore che poi viene memorizzato nella
variabile \var{errno} (sulla gestione degli errori torneremo in dettaglio in
sez.~\ref{sec:sys_errors}).

Il valore di \param{number} dipende sia dalla versione di kernel che
dall'architettura,\footnote{in genere le vecchie \textit{system call} non
  vengono eliminate e se ne aggiungono di nuove con nuovi numeri.}  ma
ciascuna \textit{system call} viene in genere identificata da una costante
nella forma \texttt{SYS\_*} dove al prefisso viene aggiunto il nome che spesso
corrisponde anche alla omonima funzione di libreria. Queste costanti sono
definite nel file \headfiled{sys/syscall.h}, ma si possono anche usare
direttamente valori numerici.


\subsection{La terminazione di un programma}
\label{sec:proc_conclusion}

Normalmente un programma conclude la sua esecuzione quando si fa ritornare la
funzione \code{main}, si usa cioè l'istruzione \instruction{return} del
linguaggio C all'interno della stessa, o se si richiede esplicitamente la
chiusura invocando direttamente la funzione \func{exit}. Queste due modalità
sono assolutamente equivalenti, dato che \func{exit} viene chiamata in maniera
trasparente anche quando \code{main} ritorna, passandogli come argomento il
valore indicato da \instruction{return}.

La funzione \funcd{exit}, che è completamente generale, essendo definita dallo
standard ANSI C, è quella che deve essere invocata per una terminazione
``\textit{normale}'', il suo prototipo è:

\begin{funcproto}{
  \fhead{unistd.h}
  \fdecl{void exit(int status)}
  \fdesc{Causa la conclusione ordinaria del programma.}
}
{La funzione non ritorna, il processo viene terminato.}
\end{funcproto}

La funzione è pensata per eseguire una conclusione pulita di un programma che
usi la libreria standard del C; essa esegue tutte le funzioni che sono state
registrate con \func{atexit} e \func{on\_exit} (vedi
sez.~\ref{sec:proc_atexit}), chiude tutti gli \textit{stream} (vedi
sez.~\ref{sec:file_stream}) effettuando il salvataggio dei dati sospesi
(chiamando \func{fclose}, vedi sez.~\ref{sec:file_fopen}), infine passa il
controllo al kernel chiamando la \textit{system call} \func{\_exit} (che
vedremo a breve) che completa la terminazione del processo.

\itindbeg{exit~status}

Il valore dell'argomento \param{status} o il valore di ritorno di \code{main}
costituisce quello che viene chiamato lo \textsl{stato di uscita}
(l'\textit{exit status}) del processo. In generale si usa questo valore per
fornire al processo padre (come vedremo in sez.~\ref{sec:proc_wait}) delle
informazioni generiche sulla riuscita o il fallimento del programma appena
terminato.

Anche se l'argomento \param{status} (ed il valore di ritorno di \code{main})
sono numeri interi di tipo \ctyp{int}, si deve tener presente che il valore
dello stato di uscita viene comunque troncato ad 8 bit, per cui deve essere
sempre compreso fra 0 e 255. Si tenga presente che se si raggiunge la fine
della funzione \code{main} senza ritornare esplicitamente si ha un valore di
uscita indefinito, è pertanto consigliabile di concludere sempre in maniera
esplicita detta funzione.

Non esiste un significato intrinseco della stato di uscita, ma una convenzione
in uso pressoché universale è quella di restituire 0 in caso di successo e 1
in caso di fallimento. Una eccezione a questa convenzione è per i programmi
che effettuano dei confronti (come \cmd{diff}), che usano 0 per indicare la
corrispondenza, 1 per indicare la non corrispondenza e 2 per indicare
l'incapacità di effettuare il confronto. Un'altra convenzione riserva i valori
da 128 a 256 per usi speciali: ad esempio 128 viene usato per indicare
l'incapacità di eseguire un altro programma in un sottoprocesso. Benché le
convenzioni citate non siano seguite universalmente è una buona idea tenerle
presenti ed adottarle a seconda dei casi.

Si tenga presente inoltre che non è una buona idea usare eventuali codici di
errore restituiti nella variabile \var{errno} (vedi sez.~\ref{sec:sys_errors})
come \textit{exit status}. In generale infatti non ci si cura del valore dello
stato di uscita di un processo se non per vedere se è diverso da zero, come
indicazione di un qualche errore.  Dato che viene troncato ad 8 bit utilizzare
un intero di valore generico può comportare il rischio, qualora si vada ad
usare un multiplo di 256, di avere uno stato di uscita uguale a zero, che
verrebbe interpretato come un successo.

Per questo motivo in \headfile{stdlib.h} sono definite, seguendo lo standard
POSIX, le due costanti \constd{EXIT\_SUCCESS} e \constd{EXIT\_FAILURE}, da
usare sempre per specificare lo stato di uscita di un processo. Su Linux, ed
in generale in qualunque sistema POSIX, ad esse sono assegnati rispettivamente
i valori 0 e 1.

\itindend{exit~status}

Una forma alternativa per effettuare una terminazione esplicita di un
programma è quella di chiamare direttamente la \textit{system call}
\funcd{\_exit},\footnote{la stessa è definita anche come \funcd{\_Exit} in
  \headfile{stdlib.h}, inoltre a partire dalla \acr{glibc} 2.3 usando questa
  funzione viene invocata \func{exit\_group} che termina tutti i
  \textit{thread} del processo e non solo quello corrente (fintanto che non si
  usano i \textit{thread}\unavref{, vedi sez.~\ref{cha:threads},} questo non
  fa nessuna differenza).} che restituisce il controllo direttamente al
kernel, concludendo immediatamente il processo, il suo prototipo è:

\begin{funcproto}{ \fhead{unistd.h} \fdecl{void \_exit(int status)}
    \fdesc{Causa la conclusione immediata del programma.}  } {La funzione non
    ritorna, il processo viene terminato.}
\end{funcproto}

La funzione termina immediatamente il processo e le eventuali funzioni
registrate con \func{atexit} e \func{on\_exit} non vengono eseguite. La
funzione chiude tutti i file descriptor appartenenti al processo, cosa che
però non comporta il salvataggio dei dati eventualmente presenti nei buffer
degli \textit{stream}, (torneremo sulle due interfacce dei file in
sez.~\ref{sec:file_unix_interface} e
sez.~\ref{sec:files_std_interface}). Infine fa sì che ogni figlio del processo
sia adottato da \cmd{init} (vedi sez.~\ref{sec:proc_termination}), manda un
segnale \signal{SIGCHLD} al processo padre (vedi
sez.~\ref{sec:sig_job_control}) e salva lo stato di uscita specificato in
\param{status} che può essere raccolto usando la funzione \func{wait} (vedi
sez.~\ref{sec:proc_wait}).

Si tenga presente infine che oltre alla conclusione ``\textsl{normale}''
appena illustrata esiste anche la possibilità di una conclusione
``\textsl{anomala}'' del programma a causa della ricezione di un segnale
(tratteremo i segnali in cap.~\ref{cha:signals}) o della chiamata alla
funzione \func{abort}; torneremo su questo in sez.~\ref{sec:proc_termination}.


\subsection{Esecuzione di funzioni preliminari all'uscita}
\label{sec:proc_atexit}

Un'esigenza comune che si incontra è quella di dover effettuare una serie di
operazioni di pulizia prima della conclusione di un programma, ad esempio
salvare dei dati, ripristinare delle impostazioni, eliminare dei file
temporanei, ecc. In genere queste operazioni vengono fatte in un'apposita
sezione del programma, ma quando si realizza una libreria diventa antipatico
dover richiedere una chiamata esplicita ad una funzione di pulizia al
programmatore che la utilizza.

È invece molto meno soggetto ad errori, e completamente trasparente
all'utente, avere la possibilità di fare effettuare automaticamente la
chiamata ad una funzione che effettui tali operazioni all'uscita dal
programma. A questo scopo lo standard ANSI C prevede la possibilità di
registrare un certo numero di funzioni che verranno eseguite all'uscita dal
programma,\footnote{nel caso di \func{atexit} lo standard POSIX.1-2001
  richiede che siano registrabili almeno \constd{ATEXIT\_MAX} funzioni (il
  valore può essere ottenuto con \func{sysconf}, vedi
  sez.~\ref{sec:sys_limits}).} sia per la chiamata ad \func{exit} che per il
ritorno di \code{main}. La prima funzione che si può utilizzare a tal fine è
\funcd{atexit}, il cui prototipo è:

\begin{funcproto}{ 
\fhead{stdlib.h} 
\fdecl{int atexit(void (*function)(void))}
\fdesc{Registra la funzione \param{function} per la chiamata all'uscita
      dal programma.}  
} 
{La funzione ritorna $0$ in caso di successo e $-1$ per un errore, \var{errno}
  non viene modificata.}
\end{funcproto}

La funzione richiede come argomento \param{function} l'indirizzo di una
opportuna funzione di pulizia da chiamare all'uscita del programma, che non
deve prendere argomenti e non deve ritornare niente. In sostanza deve la
funzione di pulizia dovrà essere definita come \code{void function(void)}.

Un'estensione di \func{atexit} è la funzione \funcd{on\_exit}, che la
\acr{glibc} include per compatibilità con SunOS ma che non è detto sia
definita su altri sistemi,\footnote{la funzione è disponibile dalla
  \acr{glibc} 2.19 definendo la macro \macro{\_DEFAULT\_SOURCE}, mentre in
  precedenza erano necessarie \macro{\_BSD\_SOURCE} o \macro{\_SVID\_SOURCE};
  non essendo prevista dallo standard POSIX è in generale preferibile evitarne
  l'uso.} il suo prototipo è:

\begin{funcproto}{ 
\fhead{stdlib.h} 
\fdecl{int on\_exit(void (*function)(int, void *), void *arg))}
\fdesc{Registra la funzione \param{function} per la chiamata all'uscita dal
  programma.} 
}
{La funzione ritorna $0$ in caso di successo e $-1$ per un errore, \var{errno}
  non viene modificata.} 
\end{funcproto}

In questo caso la funzione da chiamare all'uscita prende i due argomenti
specificati nel prototipo, un intero ed un puntatore; dovrà cioè essere
definita come \code{void function(int status, void *argp)}. Il primo argomento
sarà inizializzato allo stato di uscita con cui è stata chiamata \func{exit}
ed il secondo al puntatore \param{arg} passato come secondo argomento di
\func{on\_exit}.  Così diventa possibile passare dei dati alla funzione di
chiusura.

Nella sequenza di chiusura tutte le funzioni registrate verranno chiamate in
ordine inverso rispetto a quello di registrazione, ed una stessa funzione
registrata più volte sarà chiamata più volte. Siccome entrambe le funzioni
\func{atexit} e \func{on\_exit} fanno riferimento alla stessa lista, l'ordine
di esecuzione sarà riferito alla registrazione in quanto tale,
indipendentemente dalla funzione usata per farla.

Una volta completata l'esecuzione di tutte le funzioni registrate verranno
chiusi tutti gli \textit{stream} aperti ed infine verrà chiamata \func{\_exit}
per la terminazione del programma. Questa è la sequenza ordinaria, eseguita a
meno che una delle funzioni registrate non esegua al suo interno
\func{\_exit}, nel qual caso la terminazione del programma sarà immediata ed
anche le successive funzioni registrate non saranno invocate.

Se invece all'interno di una delle funzioni registrate si chiama un'altra
volta \func{exit} lo standard POSIX.1-2001 prescrive un comportamento
indefinito, con la possibilità (che su Linux comunque non c'è) di una
ripetizione infinita. Pertanto questa eventualità è da evitare nel modo più
assoluto. Una altro comportamento indefinito si può avere se si termina
l'esecuzione di una delle funzioni registrate con \func{longjmp} (vedi
sez.~\ref{sec:proc_longjmp}).

Si tenga presente infine che in caso di terminazione anomala di un processo
(ad esempio a causa di un segnale) nessuna delle funzioni registrate verrà
eseguita e che se invece si crea un nuovo processo con \func{fork} (vedi
sez.~\ref{sec:proc_fork}) questo manterrà tutte le funzioni già registrate.


\subsection{Un riepilogo}
\label{sec:proc_term_conclusion}

Data l'importanza dell'argomento è opportuno un piccolo riepilogo dei fatti
essenziali relativi alla esecuzione di un programma. Il primo punto da
sottolineare è che in un sistema unix-like l'unico modo in cui un programma
può essere eseguito dal kernel è attraverso la chiamata alla \textit{system
  call} \func{execve}, sia direttamente che attraverso una delle funzioni
della famiglia \func{exec} che ne semplificano l'uso (vedi
sez.~\ref{sec:proc_exec}).

Allo stesso modo l'unico modo in cui un programma può concludere
volontariamente la propria esecuzione è attraverso una chiamata alla
\textit{system call} \func{\_exit}, sia che questa venga fatta esplicitamente,
o in maniera indiretta attraverso l'uso di \func{exit} o il ritorno di
\code{main}. 

Uno schema riassuntivo che illustra le modalità con cui si avvia e conclude
normalmente un programma è riportato in fig.~\ref{fig:proc_prog_start_stop}.

\begin{figure}[htb]
  \centering
  \includegraphics[width=9cm]{img/proc_beginend}
  % \begin{tikzpicture}[>=stealth]
  %   \filldraw[fill=black!35] (-0.3,0) rectangle (12,1);
  %   \draw(5.5,0.5) node {\large{\textsf{kernel}}};

  %   \filldraw[fill=black!15] (1.5,2) rectangle (4,3);
  %   \draw (2.75,2.5) node {\texttt{ld-linux.so}};
  %   \draw [->] (2.75,1) -- (2.75,2);
  %   \draw (2.75,1.5) node [anchor=west]{\texttt{execve}};

  %   \filldraw[fill=black!15,rounded corners] (1.5,4) rectangle (4,5);
  %   \draw (2.75,4.5) node {\texttt{main}};

  %   \draw [<->, dashed] (2.75,3) -- (2.75,4);
  %   \draw [->] (1.5,4.5) -- (0.3,4.5) -- (0.3,1);
  %   \draw (0.9,4.5) node [anchor=south] {\texttt{\_exit}};

  %   \filldraw[fill=black!15,rounded corners] (1.5,6) rectangle (4,7);
  %   \draw (2.75,6.5) node {\texttt{funzione}};

  %   \draw [<->, dashed] (2.75,5) -- (2.75,6);
  %   \draw [->] (1.5,6.5) -- (0.05,6.5) -- (0.05,1);
  %   \draw (0.9,6.5) node [anchor=south] {\texttt{\_exit}};

  %   \draw (6.75,4.5) node (exit) [rectangle,fill=black!15,minimum width=2.5cm,minimum height=1cm,rounded corners, draw]{\texttt{exit}};

  %   \draw[->] (4,6.5) -- node[anchor=south west]{\texttt{exit}} (exit);
  %   \draw[->] (4,4.5) -- node[anchor=south]{\texttt{exit}} (exit);
  %   \draw[->] (exit) -- node[anchor=east]{\texttt{\_exit}}(6.75,1);

  %   \draw (10,4.5) node (exithandler1) [rectangle,fill=black!15,rounded corners, draw]{\textsf{exit handler}};
  %   \draw (10,5.5) node (exithandler2) [rectangle,fill=black!15,rounded corners, draw]{\textsf{exit handler}};
  %   \draw (10,3.5) node (stream) [rectangle,fill=black!15,rounded corners, draw]{\textsf{chiusura stream}};

  %   \draw[<->, dashed] (exithandler1) -- (exit);
  %   \draw[<->, dashed] (exithandler2) -- (exit);
  %   \draw[<->, dashed] (stream) -- (exit);
  % \end{tikzpicture}
  \caption{Schema dell'avvio e della conclusione di un programma.}
  \label{fig:proc_prog_start_stop}
\end{figure}

Si ricordi infine che un programma può anche essere interrotto dall'esterno
attraverso l'uso di un segnale (modalità di conclusione non mostrata in
fig.~\ref{fig:proc_prog_start_stop}); tratteremo nei dettagli i segnali e la
loro gestione nel capitolo \ref{cha:signals}.



\section{I processi e l'uso della memoria}
\label{sec:proc_memory}

Una delle risorse più importanti che ciascun processo ha a disposizione è la
memoria, e la gestione della memoria è appunto uno degli aspetti più complessi
di un sistema unix-like. In questa sezione, dopo una breve introduzione ai
concetti di base, esamineremo come la memoria viene vista da parte di un
programma in esecuzione, e le varie funzioni utilizzabili per la sua gestione.


\subsection{I concetti generali}
\label{sec:proc_mem_gen}

\index{memoria~virtuale|(}

Ci sono vari modi in cui i sistemi operativi organizzano la memoria, ed i
dettagli di basso livello dipendono spesso in maniera diretta
dall'architettura dell'hardware, ma quello più tipico, usato dai sistemi
unix-like come Linux è la cosiddetta \textsl{memoria virtuale} che consiste
nell'assegnare ad ogni processo uno spazio virtuale di indirizzamento lineare,
in cui gli indirizzi vanno da zero ad un qualche valore massimo.\footnote{nel
  caso di Linux fino al kernel 2.2 detto massimo era, per macchine a 32bit, di
  2Gb. Con il kernel 2.4 ed il supporto per la \textit{high-memory} il limite
  è stato esteso anche per macchine a 32 bit.}  Come accennato nel
cap.~\ref{cha:intro_unix} questo spazio di indirizzi è virtuale e non
corrisponde all'effettiva posizione dei dati nella RAM del computer. In
generale detto spazio non è neppure continuo, cioè non tutti gli indirizzi
possibili sono utilizzabili, e quelli usabili non sono necessariamente
adiacenti.

\itindbeg{huge~page}

Per la gestione da parte del kernel la memoria viene divisa in pagine di
dimensione fissa. Inizialmente queste pagine erano di 4kb sulle macchine a 32
bit e di 8kb sulle alpha. Con le versioni più recenti del kernel è possibile
anche utilizzare pagine di dimensioni maggiori (di 4Mb, dette \textit{huge
  page}), per sistemi con grandi quantitativi di memoria in cui l'uso di
pagine troppo piccole comporta una perdita di prestazioni. In alcuni sistemi
la costante \constd{PAGE\_SIZE}, definita in \headfile{limits.h}, indica la
dimensione di una pagina in byte, con Linux questo non avviene e per ottenere
questa dimensione si deve ricorrere alla funzione \func{getpagesize} (vedi
sez.~\ref{sec:sys_memory_res}).

\itindend{huge~page}
\itindbeg{page~table}

Ciascuna pagina di memoria nello spazio di indirizzi virtuale è associata ad
un supporto che può essere una pagina di memoria reale o ad un dispositivo di
stoccaggio secondario (come lo spazio disco riservato alla \textit{swap}, o i
file che contengono il codice). Per ciascun processo il kernel si cura di
mantenere un mappa di queste corrispondenze nella cosiddetta \textit{page
  table}.\footnote{questa è una semplificazione brutale, il meccanismo è molto
  più complesso; una buona trattazione di come Linux gestisce la memoria
  virtuale si trova su \cite{LinVM}.}

\itindend{page~table}

Una stessa pagina di memoria reale può fare da supporto a diverse pagine di
memoria virtuale appartenenti a processi diversi, come accade in genere per le
pagine che contengono il codice delle librerie condivise. Ad esempio il codice
della funzione \func{printf} starà su una sola pagina di memoria reale che
farà da supporto a tutte le pagine di memoria virtuale di tutti i processi che
hanno detta funzione nel loro codice.

\index{paginazione|(}

La corrispondenza fra le pagine della memoria virtuale di un processo e quelle
della memoria fisica della macchina viene gestita in maniera trasparente dal
kernel.\footnote{in genere con l'ausilio dell'hardware di gestione della
  memoria (la \textit{Memory Management Unit} del processore), con i kernel
  della serie 2.6 è comunque diventato possibile utilizzare Linux anche su
  architetture che non dispongono di una MMU.}  Poiché in genere la memoria
fisica è solo una piccola frazione della memoria virtuale, è necessario un
meccanismo che permetta di trasferire le pagine che servono dal supporto su
cui si trovano in memoria, eliminando quelle che non servono.  Questo
meccanismo è detto \textsl{paginazione} (o \textit{paging}), ed è uno dei
compiti principali del kernel.

\itindbeg{page~fault} 

Quando un processo cerca di accedere ad una pagina che non è nella memoria
reale, avviene quello che viene chiamato un \textit{page fault}; la gestione
della memoria genera un'interruzione e passa il controllo al kernel il quale
sospende il processo e si incarica di mettere in RAM la pagina richiesta,
effettuando tutte le operazioni necessarie per reperire lo spazio necessario,
per poi restituire il controllo al processo.

Dal punto di vista di un processo questo meccanismo è completamente
trasparente, e tutto avviene come se tutte le pagine fossero sempre
disponibili in memoria.  L'unica differenza avvertibile è quella dei tempi di
esecuzione, che passano dai pochi nanosecondi necessari per l'accesso in RAM
se la pagina è direttamente disponibile, a tempi estremamente più lunghi,
dovuti all'intervento del kernel, qualora sia necessario reperire pagine
riposte nella \textit{swap}.

\itindend{page~fault} 

Normalmente questo è il prezzo da pagare per avere un \textit{multitasking}
reale, ed in genere il sistema è molto efficiente in questo lavoro; quando
però ci siano esigenze specifiche di prestazioni è possibile usare delle
funzioni che permettono di bloccare il meccanismo della paginazione e
mantenere fisse delle pagine in memoria (vedi sez.~\ref{sec:proc_mem_lock}).

\index{paginazione|)}
\index{memoria~virtuale|)}


\subsection{La struttura della memoria di un processo}
\label{sec:proc_mem_layout}

\itindbeg{segment~violation}

Benché lo spazio di indirizzi virtuali copra un intervallo molto ampio, solo
una parte di essi è effettivamente allocato ed utilizzabile dal processo; il
tentativo di accedere ad un indirizzo non allocato è un tipico errore che si
commette quando si è manipolato male un puntatore e genera quella che viene
chiamata una \textit{segment violation}. Se si tenta cioè di leggere o
scrivere con un indirizzo per il quale non esiste un'associazione nella
memoria virtuale, il kernel risponde al relativo \textit{page fault} mandando
un segnale \signal{SIGSEGV} al processo, che normalmente ne causa la
terminazione immediata.

\itindend{segment~violation}

È pertanto importante capire come viene strutturata la memoria virtuale di un
processo. Essa viene divisa in \textsl{segmenti}, cioè un insieme contiguo di
indirizzi virtuali ai quali il processo può accedere.  Solitamente un
programma C viene suddiviso nei seguenti segmenti:
\index{segmento!testo|(}
\index{segmento!dati|(}
\itindbeg{heap} 
\itindbeg{stack}
\begin{enumerate}
\item Il \textsl{segmento di testo} o \textit{text segment}.  Contiene il
  codice del programma, delle funzioni di librerie da esso utilizzate, e le
  costanti.  Normalmente viene condiviso fra tutti i processi che eseguono lo
  stesso programma e nel caso delle librerie anche da processi che eseguono
  altri programmi.

  Quando l'architettura hardware lo supporta viene marcato in sola lettura per
  evitare sovrascritture accidentali (o maliziose) che ne modifichino le
  istruzioni.  Viene allocato da \func{execve} all'avvio del programma e resta
  invariato per tutto il tempo dell'esecuzione.
\index{variabili!globali|(}
\index{variabili!statiche|(}
\item Il \textsl{segmento dei dati} o \textit{data segment}. Contiene tutti i
  dati del programma, come le \textsl{variabili globali}, cioè quelle definite
  al di fuori di tutte le funzioni che compongono il programma, e le
  \textsl{variabili statiche}, cioè quelle dichiarate con l'attributo
  \direct{static},\footnote{la direttiva \direct{static} indica al compilatore
    C che una variabile così dichiarata all'interno di una funzione deve
    essere mantenuta staticamente in memoria (nel segmento dati appunto);
    questo significa che la variabile verrà inizializzata una sola volta alla
    prima invocazione della funzione e che il suo valore sarà mantenuto fra
    diverse esecuzioni della funzione stessa, la differenza con una variabile
    globale è che essa può essere vista solo all'interno della funzione in cui
    è dichiarata.} e la memoria allocata dinamicamente. Di norma è diviso in
  tre parti:
  \begin{itemize}
  \item Il segmento dei dati inizializzati, che contiene le variabili il cui
    valore è stato assegnato esplicitamente. Ad esempio se si definisce:
    \includecodesnip{listati/pi.c}
    questo valore sarà immagazzinato in questo segmento. La memoria di questo
    segmento viene preallocata all'avvio del programma e inizializzata ai valori
    specificati.
  \item Il segmento dei dati non inizializzati, che contiene le variabili il
    cui valore non è stato assegnato esplicitamente. Ad esempio se si
    definisce:
    \includecodesnip{listati/vect.c}
    questo vettore sarà immagazzinato in questo segmento. Anch'esso viene
    allocato all'avvio, e tutte le variabili vengono inizializzate a zero (ed
    i puntatori a \val{NULL}).\footnote{si ricordi che questo vale solo per le
      variabili che vanno nel segmento dati, e non è affatto vero in
      generale.}  Storicamente questa seconda parte del segmento dati viene
    chiamata \itindex{Block~Started~by~Symbol~(BSS)} BSS (da \textit{Block
      Started by Symbol}). La sua dimensione è fissa.
    \index{variabili!globali|)} \index{variabili!statiche|)}
  \item Lo \textit{heap}, detto anche \textit{free store}. Tecnicamente lo si
    può considerare l'estensione del segmento dei dati non inizializzati, a
    cui di solito è posto giusto di seguito. Questo è il segmento che viene
    utilizzato per l'allocazione dinamica della memoria.  Lo \textit{heap} può
    essere ridimensionato allargandolo e restringendolo per allocare e
    disallocare la memoria dinamica con le apposite funzioni (vedi
    sez.~\ref{sec:proc_mem_alloc}), ma il suo limite inferiore, quello
    adiacente al segmento dei dati non inizializzati, ha una posizione fissa.
  \end{itemize}
\item Il segmento di \textit{stack}, che contiene quello che viene chiamato lo
  ``\textit{stack}'' del programma.  Tutte le volte che si effettua una
  chiamata ad una funzione è qui che viene salvato l'indirizzo di ritorno e le
  informazioni dello stato del chiamante (come il contenuto di alcuni registri
  della CPU), poi la funzione chiamata alloca qui lo spazio per le sue
  variabili locali. Tutti questi dati vengono \textit{impilati} (da questo
  viene il nome \textit{stack}) in sequenza uno sull'altro; in questo modo le
  funzioni possono essere chiamate ricorsivamente. Al ritorno della funzione
  lo spazio è automaticamente rilasciato e ``\textsl{ripulito}''.\footnote{il
    compilatore si incarica di generare automaticamente il codice necessario,
    seguendo quella che viene chiamata una \textit{calling convention}; quella
    standard usata con il C ed il C++ è detta \textit{cdecl} e prevede che gli
    argomenti siano caricati nello \textit{stack} dal chiamante da destra a
    sinistra, e che sia il chiamante stesso ad eseguire la ripulitura dello
    \textit{stack} al ritorno della funzione, se ne possono però utilizzare di
    alternative (ad esempio nel Pascal gli argomenti sono inseriti da sinistra
    a destra ed è compito del chiamato ripulire lo \textit{stack}), in genere
    non ci si deve preoccupare di questo fintanto che non si mescolano
    funzioni scritte con linguaggi diversi.}

  La dimensione di questo segmento aumenta seguendo la crescita dello
  \textit{stack} del programma, ma non viene ridotta quando quest'ultimo si
  restringe.
\end{enumerate}

\begin{figure}[htb]
  \centering
  \includegraphics[height=10cm]{img/memory_layout}
  % \begin{tikzpicture}
  % \draw (0,0) rectangle (4,1);
  % \draw (2,0.5) node {\textit{text}};
  % \draw (0,1) rectangle (4,2.5);
  % \draw (2,1.75) node {dati inizializzati};
  % \draw (0,2.5) rectangle (4,5);
  % \draw (2,3.75) node {dati non inizializzati};
  % \draw (0,5) rectangle (4,9);
  % \draw[dashed] (0,6) -- (4,6);
  % \draw[dashed] (0,8) -- (4,8);
  % \draw (2,5.5) node {\textit{heap}};
  % \draw (2,8.5) node {\textit{stack}};
  % \draw [->] (2,6) -- (2,6.5);
  % \draw [->] (2,8) -- (2,7.5);
  % \draw (0,9) rectangle (4,10);
  % \draw (2,9.5) node {\textit{environment}};
  % \draw (4,0) node [anchor=west] {\texttt{0x08000000}};
  % \draw (4,5) node [anchor=west] {\texttt{0x08xxxxxx}};
  % \draw (4,9) node [anchor=west] {\texttt{0xC0000000}};
  % \end{tikzpicture} 
  \caption{Disposizione tipica dei segmenti di memoria di un processo.}
  \label{fig:proc_mem_layout}
\end{figure}

Una disposizione tipica dei vari segmenti (testo, dati inizializzati e non
inizializzati, \textit{heap}, \textit{stack}, ecc.) è riportata in
fig.~\ref{fig:proc_mem_layout}. Si noti come in figura sia indicata una
ulteriore regione, marcata \textit{environment}, che è quella che contiene i
dati relativi alle variabili di ambiente passate al programma al suo avvio
(torneremo su questo argomento in sez.~\ref{sec:proc_environ}).

Usando il comando \cmd{size} su un programma se ne può stampare le dimensioni
dei segmenti di testo e di dati (solo però per i dati inizializzati ed il BSS,
dato che lo \textit{heap} ha una dimensione dinamica). Si tenga presente
comunque che il BSS, contrariamente al segmento dei dati inizializzati, non è
mai salvato sul file che contiene l'eseguibile, dato che viene sempre
inizializzato a zero al caricamento del programma.

\index{segmento!testo|)}
\index{segmento!dati|)}
\itindend{heap} 
\itindend{stack}


\subsection{Allocazione della memoria per i programmi C}
\label{sec:proc_mem_alloc}

Il C supporta direttamente, come linguaggio di programmazione, soltanto due
modalità di allocazione della memoria: l'\textsl{allocazione statica} e
l'\textsl{allocazione automatica}.

L'\textsl{allocazione statica} è quella con cui sono memorizzate le variabili
globali e le variabili statiche, cioè le variabili il cui valore deve essere
mantenuto per tutta la durata del programma. Come accennato queste variabili
vengono allocate nel segmento dei dati all'avvio del programma come parte
delle operazioni svolte da \func{exec}, e lo spazio da loro occupato non viene
liberato fino alla sua conclusione.

\index{variabili!automatiche|(}

L'\textsl{allocazione automatica} è quella che avviene per gli argomenti di
una funzione e per le sue variabili locali, quelle che vengono definite
all'interno della funzione che esistono solo per la durata della sua
esecuzione e che per questo vengono anche dette \textsl{variabili
  automatiche}.  Lo spazio per queste variabili viene allocato nello
\textit{stack} quando viene eseguita la funzione e liberato quando si esce
dalla medesima.

\index{variabili!automatiche|)}

Esiste però un terzo tipo di allocazione, l'\textsl{allocazione dinamica}
della memoria, che non è prevista direttamente all'interno del linguaggio C,
ma che è necessaria quando il quantitativo di memoria che serve è
determinabile solo durante il corso dell'esecuzione del programma. Il C non
consente di usare variabili allocate dinamicamente, non è possibile cioè
definire in fase di programmazione una variabile le cui dimensioni possano
essere modificate durante l'esecuzione del programma. Per questo la libreria
standard del C fornisce una opportuna serie di funzioni per eseguire
l'allocazione dinamica di memoria, che come accennato avviene nello
\textit{heap}.

Le variabili il cui contenuto è allocato in questo modo non potranno essere
usate direttamente come le altre (quelle nello \textit{stack}), ma l'accesso
sarà possibile solo in maniera indiretta, attraverso i puntatori alla memoria
loro riservata che si sono ottenuti dalle funzioni di allocazione.

Le funzioni previste dallo standard ANSI C per la gestione della memoria sono
quattro: \func{malloc}, \func{calloc}, \func{realloc} e \func{free}. Le prime
due, \funcd{malloc} e \funcd{calloc} allocano nuovo spazio di memoria; i
rispettivi prototipi sono:

\begin{funcproto}{ 
\fhead{stdlib.h} 
\fdecl{void *calloc(size\_t nmemb, size\_t size)}
\fdesc{Alloca un'area di memoria inizializzata a 0.}  
\fdecl{void *malloc(size\_t size)}
\fdesc{Alloca un'area di memoria non inizializzata.}  
}
{Entrambe le funzioni restituiscono il puntatore alla zona di memoria allocata
in caso di successo e \val{NULL} in caso di fallimento, nel qual caso
  \var{errno} assumerà il valore \errcode{ENOMEM}.}
\end{funcproto}

In genere si usano \func{malloc} e \func{calloc} per allocare dinamicamente
un'area di memoria.\footnote{queste funzioni presentano un comportamento
  diverso fra la \acr{glibc} e la \acr{uClib} quando il valore di \param{size}
  è nullo.  Nel primo caso viene comunque restituito un puntatore valido,
  anche se non è chiaro a cosa esso possa fare riferimento, nel secondo caso
  viene restituito \val{NULL}. Il comportamento è analogo con
  \code{realloc(NULL, 0)}.}  Dato che i puntatori ritornati sono di tipo
generico non è necessario effettuare un cast per assegnarli a puntatori al
tipo di variabile per la quale si effettua l'allocazione, inoltre le funzioni
garantiscono che i puntatori siano allineati correttamente per tutti i tipi di
dati; ad esempio sulle macchine a 32 bit in genere sono allineati a multipli
di 4 byte e sulle macchine a 64 bit a multipli di 8 byte.

Nel caso di \func{calloc} l'area di memoria viene allocata nello \textit{heap}
come un vettore di \param{nmemb} membri di \param{size} byte di dimensione, e
preventivamente inizializzata a zero, nel caso di \func{malloc} invece vengono
semplicemente allocati \param{size} byte e l'area di memoria non viene
inizializzata.

Una volta che non sia più necessaria la memoria allocata dinamicamente deve
essere esplicitamente rilasciata usando la funzione \funcd{free},\footnote{le
  glibc provvedono anche una funzione \funcm{cfree} definita per compatibilità
  con SunOS, che è deprecata.} il suo prototipo è:

\begin{funcproto}{ 
\fhead{stdlib.h} 
\fdecl{void free(void *ptr)}
\fdesc{Disalloca un'area di memoria precedentemente allocata.}  
}
{La funzione non ritorna nulla e non riporta errori.}
\end{funcproto}

Questa funzione vuole come argomento \var{ptr} il puntatore restituito da una
precedente chiamata ad una qualunque delle funzioni di allocazione che non sia
già stato liberato da un'altra chiamata a \func{free}. Se il valore di
\param{ptr} è \val{NULL} la funzione non fa niente, mentre se l'area di
memoria era già stata liberata da una precedente chiamata il comportamento
della funzione è dichiarato indefinito, ma in genere comporta la corruzione
dei dati di gestione dell'allocazione, che può dar luogo a problemi gravi, ad
esempio un \textit{segmentation fault} in una successiva chiamata di una di
queste funzioni.

\itindbeg{double~free}

Dato che questo errore, chiamato in gergo \textit{double free}, è abbastanza
frequente, specie quando si manipolano vettori di puntatori, e dato che le
conseguenze possono essere pesanti ed inaspettate, si suggerisce come
soluzione precauzionale di assegnare sempre a \val{NULL} ogni puntatore su cui
sia stata eseguita \func{free} immediatamente dopo l'esecuzione della
funzione. In questo modo, dato che con un puntatore nullo \func{free} non
esegue nessuna operazione, si evitano i problemi del \textit{double free}.

\itindend{double~free}

Infine la funzione \funcd{realloc} consente di modificare, in genere di
aumentare, la dimensione di un'area di memoria precedentemente allocata; il
suo prototipo è:

\begin{funcproto}{ 
\fhead{stdlib.h} 
\fdecl{void *realloc(void *ptr, size\_t size)}
\fdesc{Cambia la dimensione di un'area di memoria precedentemente allocata.}
}  {La funzione ritorna il puntatore alla zona di memoria allocata in caso
  di successo e \val{NULL} per un errore, nel qual caso \var{errno}
  assumerà il valore \errcode{ENOMEM}.}
\end{funcproto}

La funzione vuole come primo argomento il puntatore restituito da una
precedente chiamata a \func{malloc} o \func{calloc} e come secondo argomento
la nuova dimensione (in byte) che si intende ottenere. Se si passa
per \param{ptr} il valore \val{NULL} allora la funzione si comporta come
\func{malloc}.\footnote{questo è vero per Linux e l'implementazione secondo lo
  standard ANSI C, ma non è vero per alcune vecchie implementazioni, inoltre
  alcune versioni delle librerie del C consentivano di usare \func{realloc}
  anche per un puntatore liberato con \func{free} purché non ci fossero state
  nel frattempo altre chiamate a funzioni di allocazione, questa funzionalità
  è totalmente deprecata e non è consentita sotto Linux.}

La funzione si usa ad esempio quando si deve far crescere la dimensione di un
vettore. In questo caso se è disponibile dello spazio adiacente al precedente
la funzione lo utilizza, altrimenti rialloca altrove un blocco della
dimensione voluta, copiandoci automaticamente il contenuto; lo spazio aggiunto
non viene inizializzato. Se la funzione fallisce l'area di memoria originale
non viene assolutamente toccata.

Si deve sempre avere ben presente il fatto che il blocco di memoria restituito
da \func{realloc} può non essere un'estensione di quello che gli si è passato
in ingresso; per questo si dovrà \emph{sempre} eseguire la riassegnazione di
\param{ptr} al valore di ritorno della funzione, e reinizializzare o provvedere
ad un adeguato aggiornamento di tutti gli altri puntatori all'interno del
blocco di dati ridimensionato.

La \acr{glibc} ha un'implementazione delle funzioni di allocazione che è
controllabile dall'utente attraverso alcune variabili di ambiente (vedi
sez.~\ref{sec:proc_environ}), in particolare diventa possibile tracciare
questo tipo di errori usando la variabile di ambiente \envvar{MALLOC\_CHECK\_}
che quando viene definita mette in uso una versione meno efficiente delle
funzioni suddette, che però è più tollerante nei confronti di piccoli errori
come quello dei \textit{double free} o i \textit{buffer overrun} di un
byte.\footnote{uno degli errori più comuni, causato ad esempio dalla scrittura
  di una stringa di dimensione pari a quella del buffer, in cui ci si
  dimentica dello zero di terminazione finale.}  In particolare:
\begin{itemize*}
\item se la variabile è posta a $0$ gli errori vengono ignorati;
\item se la variabile è posta a $1$ viene stampato un avviso sullo
  \textit{standard error} (vedi sez.~\ref{sec:file_fd});
\item se la variabile è posta a $2$ viene chiamata la funzione \func{abort}
  (vedi sez.~\ref{sec:sig_alarm_abort}), che in genere causa l'immediata
  terminazione del programma;
\item se la variabile è posta a $3$ viene stampato l'avviso e chiamata
  \func{abort}. 
\end{itemize*}

\itindbeg{memory~leak}

L'errore di programmazione più comune e più difficile da risolvere che si
incontra con le funzioni di allocazione è quando non viene opportunamente
liberata la memoria non più utilizzata, quello che in inglese viene chiamato
\textit{memory leak}, cioè una \textsl{perdita di memoria}.

Un caso tipico che illustra il problema è quello in cui in una propria
funzione si alloca della memoria per uso locale senza liberarla prima di
uscire. La memoria resta così allocata fino alla terminazione del processo.
Chiamate ripetute alla stessa funzione continueranno ad effettuare altre
allocazioni, che si accumuleranno causando a lungo andare un esaurimento della
memoria disponibile e la probabile impossibilità di proseguire l'esecuzione
del programma.

Il problema è che l'esaurimento della memoria può avvenire in qualunque
momento, in corrispondenza ad una qualunque chiamata di \func{malloc} che può
essere in una sezione del codice che non ha alcuna relazione con la funzione
che contiene l'errore. Per questo motivo è sempre molto difficile trovare un
\textit{memory leak}.  In C e C++ il problema è particolarmente sentito. In
C++, per mezzo della programmazione ad oggetti, il problema dei \textit{memory
  leak} si può notevolmente ridimensionare attraverso l'uso accurato di
appositi oggetti come gli \textit{smartpointers}.  Questo però in genere va a
scapito delle prestazioni dell'applicazione in esecuzione.

% TODO decidere cosa fare di questo che segue In altri linguaggi come il java
% e recentemente il C\# il problema non si pone nemmeno perché la gestione
% della memoria viene fatta totalmente in maniera automatica, ovvero il
% programmatore non deve minimamente preoccuparsi di liberare la memoria
% allocata precedentemente quando non serve più, poiché l'infrastruttura del
% linguaggio gestisce automaticamente la cosiddetta
% \itindex{garbage~collection} \textit{garbage collection}. In tal caso,
% attraverso meccanismi simili a quelli del \textit{reference counting},
% quando una zona di memoria precedentemente allocata non è più riferita da
% nessuna parte del codice in esecuzione, può essere deallocata
% automaticamente in qualunque momento dall'infrastruttura.

% Anche questo va a scapito delle prestazioni dell'applicazione in esecuzione
% (inoltre le applicazioni sviluppate con tali linguaggi di solito non sono
% eseguibili compilati, come avviene invece per il C ed il C++, ed è necessaria
% la presenza di una infrastruttura per la loro interpretazione e pertanto hanno
% di per sé delle prestazioni più scadenti rispetto alle stesse applicazioni
% compilate direttamente).  Questo comporta però il problema della non
% predicibilità del momento in cui viene deallocata la memoria precedentemente
% allocata da un oggetto.

Per limitare l'impatto di questi problemi, e semplificare la ricerca di
eventuali errori, l'implementazione delle funzioni di allocazione nella
\acr{glibc} mette a disposizione una serie di funzionalità che permettono di
tracciare le allocazioni e le disallocazioni, e definisce anche una serie di
possibili \textit{hook} (\textsl{ganci}) che permettono di sostituire alle
funzioni di libreria una propria versione (che può essere più o meno
specializzata per il debugging). Esistono varie librerie che forniscono dei
sostituti opportuni delle funzioni di allocazione in grado, senza neanche
ricompilare il programma,\footnote{esempi sono \textit{Dmalloc}
  \url{http://dmalloc.com/} di Gray Watson ed \textit{Electric Fence} di Bruce
  Perens.} di eseguire diagnostiche anche molto complesse riguardo
l'allocazione della memoria. Vedremo alcune delle funzionalità di ausilio
presenti nella \acr{glibc} in sez.~\ref{sec:proc_memory_adv_management}.

\itindend{memory~leak}

Una possibile alternativa all'uso di \func{malloc}, per evitare di soffrire
dei problemi di \textit{memory leak} descritti in precedenza, è di allocare la
memoria nel segmento di \textit{stack} della funzione corrente invece che
nello \textit{heap}. Per farlo si può usare la funzione \funcd{alloca}, la cui
sintassi è identica a quella di \func{malloc}; il suo prototipo è:

\begin{funcproto}{ 
\fhead{stdlib.h} 
\fdecl{void *alloca(size\_t size)}
\fdesc{Alloca un'area di memoria nello \textit{stack}.} 
}
{La funzione ritorna il puntatore alla zona di memoria allocata, in caso
  di errore il comportamento è indefinito.}
\end{funcproto}

La funzione alloca la quantità di memoria (non inizializzata) richiesta
dall'argomento \param{size} nel segmento di \textit{stack} della funzione
chiamante. Con questa funzione non è più necessario liberare la memoria
allocata, e quindi non esiste un analogo della \func{free}, in quanto essa
viene rilasciata automaticamente al ritorno della funzione.

Come è evidente questa funzione ha alcuni vantaggi interessanti, anzitutto
permette di evitare alla radice i problemi di \textit{memory leak}, dato che
non serve più la deallocazione esplicita; inoltre la deallocazione automatica
funziona anche quando si usa \func{longjmp} per uscire da una subroutine con
un salto non locale da una funzione (vedi sez.~\ref{sec:proc_longjmp}).  Un
altro vantaggio è che in Linux la funzione è molto più veloce di \func{malloc}
e non viene sprecato spazio, infatti non è necessario gestire un pool di
memoria da riservare e si evitano così anche i problemi di frammentazione di
quest'ultimo, che comportano inefficienze sia nell'allocazione della memoria
che nell'esecuzione dell'allocazione.

Gli svantaggi sono che questa funzione non è disponibile su tutti gli Unix, e
non è inserita né nello standard POSIX né in SUSv3 (ma è presente in BSD), il
suo utilizzo quindi limita la portabilità dei programmi. Inoltre la funzione
non può essere usata nella lista degli argomenti di una funzione, perché lo
spazio verrebbe allocato nel mezzo degli stessi. Inoltre non è chiaramente
possibile usare \func{alloca} per allocare memoria che deve poi essere usata
anche al di fuori della funzione in cui essa viene chiamata, dato che
all'uscita dalla funzione lo spazio allocato diventerebbe libero, e potrebbe
essere sovrascritto all'invocazione di nuove funzioni.  Questo è lo stesso
problema che si può avere con le variabili automatiche, su cui torneremo in
sez.~\ref{sec:proc_var_passing}.

Infine non esiste un modo di sapere se l'allocazione ha avuto successo, la
funzione infatti viene realizzata inserendo del codice \textit{inline} nel
programma\footnote{questo comporta anche il fatto che non è possibile
  sostituirla con una propria versione o modificarne il comportamento
  collegando il proprio programma con un'altra libreria.} che si limita a
modificare il puntatore nello \textit{stack} e non c'è modo di sapere se se ne
sono superate le dimensioni, per cui in caso di fallimento nell'allocazione il
comportamento del programma può risultare indefinito, dando luogo ad una
\textit{segment violation} la prima volta che si cerchi di accedere alla
memoria non effettivamente disponibile.

\index{segmento!dati|(}
\itindbeg{heap} 

Le due funzioni seguenti vengono utilizzate soltanto quando è necessario
effettuare direttamente la gestione della memoria associata allo spazio dati
di un processo,\footnote{le due funzioni sono state definite con BSD 4.3, sono
  marcate obsolete in SUSv2 e non fanno parte delle librerie standard del C e
  mentre sono state esplicitamente rimosse dallo standard POSIX.1-2001.} per
poterle utilizzare è necessario definire una della macro di funzionalità (vedi
sez.~\ref{sec:intro_gcc_glibc_std}) fra \macro{\_BSD\_SOURCE},
\macro{\_SVID\_SOURCE} e \macro{\_XOPEN\_SOURCE} (ad un valore maggiore o
uguale di 500). La prima funzione è \funcd{brk}, ed il suo prototipo è:

\begin{funcproto}{ 
\fhead{unistd.h} 
\fdecl{int brk(void *addr)}
\fdesc{Sposta la fine del segmento dati del processo.} 
}
{La funzione ritorna $0$ in caso di successo e $-1$ per un errore,
  nel qual caso \var{errno} assumerà il valore \errcode{ENOMEM}.}
\end{funcproto}

La funzione è un'interfaccia all'omonima \textit{system call} ed imposta
l'indirizzo finale del segmento dati di un processo (più precisamente dello
\textit{heap}) all'indirizzo specificato da \param{addr}. Quest'ultimo deve
essere un valore ragionevole e la dimensione totale non deve comunque eccedere
un eventuale limite (vedi sez.~\ref{sec:sys_resource_limit}) sulle dimensioni
massime del segmento dati del processo.

Il valore di ritorno della funzione fa riferimento alla versione fornita dalla
\acr{glibc}, in realtà in Linux la \textit{system call} corrispondente
restituisce come valore di ritorno il nuovo valore della fine del segmento
dati in caso di successo e quello corrente in caso di fallimento, è la
funzione di interfaccia usata dalla \acr{glibc} che fornisce i valori di
ritorno appena descritti; se si usano librerie diverse questo potrebbe non
accadere.

Una seconda funzione per la manipolazione diretta delle dimensioni del
segmento dati\footnote{in questo caso si tratta soltanto di una funzione di
  libreria, anche se basata sulla stessa \textit{system call}.} è
\funcd{sbrk}, ed il suo prototipo è:

\begin{funcproto}{ 
\fhead{unistd.h} 
\fdecl{void *sbrk(intptr\_t increment)}
\fdesc{Incrementa la dimensione del segmento dati del processo.} 
}
{La funzione ritorna il puntatore all'inizio della nuova zona di memoria
  allocata in caso di successo e \val{NULL} per un errore, nel qual
  caso \var{errno} assumerà il valore \errcode{ENOMEM}.}
\end{funcproto}

La funzione incrementa la dimensione dello \textit{heap} di un programma del
valore indicato dall'argomento \param{increment}, restituendo il nuovo
indirizzo finale dello stesso.  L'argomento è definito come di tipo
\typed{intptr\_t}, ma a seconda della versione delle librerie e del sistema
può essere indicato con una serie di tipi equivalenti come \type{ptrdiff\_t},
\type{ssize\_t}, \ctyp{int}. Se invocata con un valore nullo la funzione
permette di ottenere l'attuale posizione della fine del segmento dati.

Queste due funzioni sono state deliberatamente escluse dallo standard POSIX.1
dato che per i normali programmi è sempre opportuno usare le funzioni di
allocazione standard descritte in precedenza, a meno di non voler realizzare
per proprio conto un diverso meccanismo di gestione della memoria del segmento
dati.
\itindend{heap} 
\index{segmento!dati|)}


\subsection{Il controllo della memoria virtuale}  
\label{sec:proc_mem_lock}

\index{memoria~virtuale|(}

Come spiegato in sez.~\ref{sec:proc_mem_gen} il kernel gestisce la memoria
virtuale in maniera trasparente ai processi, decidendo quando rimuovere pagine
dalla memoria per metterle nell'area di \textit{swap}, sulla base
dell'utilizzo corrente da parte dei vari processi.

Nell'uso comune un processo non deve preoccuparsi di tutto ciò, in quanto il
meccanismo della paginazione riporta in RAM, ed in maniera trasparente, tutte
le pagine che gli occorrono; esistono però esigenze particolari in cui non si
vuole che questo meccanismo si attivi. In generale i motivi per cui si possono
avere di queste necessità sono due:
\begin{itemize*}
\item \textsl{La velocità}. Il processo della paginazione è trasparente solo
  se il programma in esecuzione non è sensibile al tempo che occorre a
  riportare la pagina in memoria; per questo motivo processi critici che hanno
  esigenze di tempo reale o tolleranze critiche nelle risposte (ad esempio
  processi che trattano campionamenti sonori) possono non essere in grado di
  sopportare le variazioni della velocità di accesso dovuta alla paginazione.
  
  In certi casi poi un programmatore può conoscere meglio dell'algoritmo di
  allocazione delle pagine le esigenze specifiche del suo programma e decidere
  quali pagine di memoria è opportuno che restino in memoria per un aumento
  delle prestazioni. In genere queste sono esigenze particolari e richiedono
  anche un aumento delle priorità in esecuzione del processo (vedi
  sez.~\ref{sec:proc_real_time}).
  
\item \textsl{La sicurezza}. Se si hanno password o chiavi segrete in chiaro
  in memoria queste possono essere portate su disco dal meccanismo della
  paginazione. Questo rende più lungo il periodo di tempo in cui detti segreti
  sono presenti in chiaro e più complessa la loro cancellazione: un processo
  infatti può cancellare la memoria su cui scrive le sue variabili, ma non può
  toccare lo spazio disco su cui una pagina di memoria può essere stata
  salvata. Per questo motivo di solito i programmi di crittografia richiedono
  il blocco di alcune pagine di memoria.
\end{itemize*}

Per ottenere informazioni sulle modalità in cui un programma sta usando la
memoria virtuale è disponibile una apposita funzione di sistema,
\funcd{mincore}, che però non è standardizzata da POSIX e pertanto non è
disponibile su tutte le versioni di kernel unix-like;\footnote{nel caso di
  Linux devono essere comunque definite le macro \macro{\_BSD\_SOURCE} e
  \macro{\_SVID\_SOURCE} o \macro{\_DEFAULT\_SOURCE}.}  il suo prototipo è:

\begin{funcproto}{
\fhead{unistd.h}
\fhead{sys/mman.h}
\fdecl{int mincore(void *addr, size\_t length, unsigned char *vec)}
\fdesc{Ritorna lo stato delle pagine di memoria occupate da un processo.}
}
{La funzione ritorna $0$ in caso di successo e $-1$ per un errore, nel qual
caso \var{errno} assumerà uno dei valori:
\begin{errlist}
   \item[\errcode{EAGAIN}] il kernel è temporaneamente non in grado di fornire
     una risposta.
   \item[\errcode{EFAULT}] \param{vec} punta ad un indirizzo non valido.
   \item[\errcode{EINVAL}] \param{addr} non è un multiplo delle dimensioni di
     una pagina.
   \item[\errcode{ENOMEM}] o \param{addr}$+$\param{length} eccede la dimensione
     della memoria usata dal processo o l'intervallo di indirizzi specificato
     non è mappato.
\end{errlist}}
\end{funcproto}

La funzione permette di ottenere le informazioni sullo stato della mappatura
della memoria per il processo chiamante, specificando l'intervallo da
esaminare con l'indirizzo iniziale, indicato con l'argomento \param{addr}, e
la lunghezza, indicata con l'argomento \param{length}. L'indirizzo iniziale
deve essere un multiplo delle dimensioni di una pagina, mentre la lunghezza
può essere qualunque, fintanto che si resta nello spazio di indirizzi del
processo,\footnote{in caso contrario si avrà un errore di \errcode{ENOMEM};
  fino al kernel 2.6.11 in questo caso veniva invece restituito
  \errcode{EINVAL}, in considerazione che il caso più comune in cui si
  verifica questo errore è quando si usa per sbaglio un valore negativo
  di \param{length}, che nel caso verrebbe interpretato come un intero
  positivo di grandi dimensioni.}  ma il risultato verrà comunque fornito per
l'intervallo compreso fino al multiplo successivo.

I risultati della funzione vengono forniti nel vettore puntato da \param{vec},
che deve essere allocato preventivamente e deve essere di dimensione
sufficiente a contenere tanti byte quante sono le pagine contenute
nell'intervallo di indirizzi specificato, la dimensione cioè deve essere
almeno pari a \code{(length+PAGE\_SIZE-1)/PAGE\_SIZE}.  Al ritorno della
funzione il bit meno significativo di ciascun byte del vettore sarà acceso se
la pagina di memoria corrispondente è al momento residente in memoria, o
cancellato altrimenti. Il comportamento sugli altri bit è indefinito, essendo
questi al momento riservati per usi futuri. Per questo motivo in genere è
comunque opportuno inizializzare a zero il contenuto del vettore, così che le
pagine attualmente residenti in memoria saranno indicate da un valore non
nullo del byte corrispondente.

Dato che lo stato della memoria di un processo può cambiare continuamente, il
risultato di \func{mincore} è assolutamente provvisorio e lo stato delle
pagine potrebbe essere già cambiato al ritorno stesso della funzione, a meno
che, come vedremo ora, non si sia attivato il meccanismo che forza il
mantenimento di una pagina sulla memoria.  

\itindbeg{memory~locking}

Il meccanismo che previene la paginazione di parte della memoria virtuale di
un processo è chiamato \textit{memory locking} (o \textsl{blocco della
  memoria}). Il blocco è sempre associato alle pagine della memoria virtuale
del processo, e non al segmento reale di RAM su cui essa viene mantenuta.  La
regola è che se un segmento di RAM fa da supporto ad almeno una pagina
bloccata allora esso viene escluso dal meccanismo della paginazione. I blocchi
non si accumulano, se si blocca due volte la stessa pagina non è necessario
sbloccarla due volte, una pagina o è bloccata oppure no.

Il \textit{memory lock} persiste fintanto che il processo che detiene la
memoria bloccata non la sblocca. Chiaramente la terminazione del processo
comporta anche la fine dell'uso della sua memoria virtuale, e quindi anche di
tutti i suoi \textit{memory lock}.  Inoltre i \textit{memory lock} non sono
ereditati dai processi figli, ma siccome Linux usa il \textit{copy on write}
(vedi sez.~\ref{sec:proc_fork}) gli indirizzi virtuali del figlio sono
mantenuti sullo stesso segmento di RAM del padre, e quindi fintanto che un
figlio non scrive su un segmento bloccato, può usufruire del \textit{memory
  lock} del padre. Infine i \textit{memory lock} vengono automaticamente
rimossi se si pone in esecuzione un altro programma con \func{exec} (vedi
sez.~\ref{sec:proc_exec}).

Il sistema pone dei limiti all'ammontare di memoria di un processo che può
essere bloccata e al totale di memoria fisica che si può dedicare a questo, lo
standard POSIX.1 richiede che sia definita in \headfile{unistd.h} la macro
\macrod{\_POSIX\_MEMLOCK\_RANGE} per indicare la capacità di eseguire il
\textit{memory locking}.

Siccome la richiesta di un \textit{memory lock} da parte di un processo riduce
la memoria fisica disponibile nel sistema per gli altri processi, questo ha un
evidente impatto su tutti gli altri processi, per cui fino al kernel 2.6.9
solo un processo dotato di privilegi amministrativi (la \textit{capability}
\const{CAP\_IPC\_LOCK}, vedi sez.~\ref{sec:proc_capabilities}) aveva la
capacità di bloccare una pagina di memoria.

A partire dal kernel 2.6.9 anche un processo normale può bloccare la propria
memoria\footnote{la funzionalità è stata introdotta per non essere costretti a
  dare privilegi eccessivi a programmi di crittografia, che necessitano di
  questa funzionalità, ma che devono essere usati da utenti normali.} ma
mentre un processo privilegiato non ha limiti sulla quantità di memoria che
può bloccare, un processo normale è soggetto al limite della risorsa
\const{RLIMIT\_MEMLOCK} (vedi sez.~\ref{sec:sys_resource_limit}). In generale
poi ogni processo può sbloccare le pagine relative alla propria memoria, se
però diversi processi bloccano la stessa pagina questa resterà bloccata
fintanto che ci sarà almeno un processo che la blocca.

Le funzioni di sistema per bloccare e sbloccare la paginazione di singole
sezioni di memoria sono rispettivamente \funcd{mlock} e \funcd{munlock}; i
loro prototipi sono:

\begin{funcproto}{
  \fhead{sys/mman.h} 
  \fdecl{int mlock(const void *addr, size\_t len)}
  \fdesc{Blocca la paginazione su un intervallo di memoria.}

  \fdecl{int munlock(const void *addr, size\_t len)}
  \fdesc{Rimuove il blocco della paginazione su un intervallo di memoria.}
  }
{Entrambe le funzioni ritornano $0$ in caso di successo e $-1$ in caso di
  errore, nel qual caso \var{errno} assumerà uno dei valori:
  \begin{errlist}
  \item[\errcode{EAGAIN}] una parte o tutto l'intervallo richiesto non può
    essere bloccato per una mancanza temporanea di risorse.
  \item[\errcode{EINVAL}] \param{len} non è un valore positivo o la somma con
    \param{addr} causa un overflow.
  \item[\errcode{ENOMEM}] alcuni indirizzi dell’intervallo specificato non
    corrispondono allo spazio di indirizzi del processo o con \func{mlock} si
    è superato il limite di \const{RLIMIT\_MEMLOCK} per un processo non
    privilegiato (solo per kernel a partire dal 2.6.9) o si è superato il
    limite di regioni di memoria con attributi diversi.
  \item[\errcode{EPERM}] il processo non è privilegiato (per kernel precedenti
    il 2.6.9) o si ha un limite nullo per \const{RLIMIT\_MEMLOCK} e
    il processo non è privilegiato (per kernel a partire dal 2.6.9).
  \end{errlist}}
\end{funcproto}

Le due funzioni permettono rispettivamente di bloccare e sbloccare la
paginazione per l'intervallo di memoria iniziante all'indirizzo \param{addr} e
lungo \param{len} byte.  Al ritorno di \func{mlock} tutte le pagine che
contengono una parte dell'intervallo bloccato sono garantite essere in RAM e
vi verranno mantenute per tutta la durata del blocco. Con kernel diversi da
Linux si può ottenere un errore di \errcode{EINVAL} se \param{addr} non è un
multiplo della dimensione delle pagine di memoria, pertanto se si ha a cuore
la portabilità si deve avere cura di allinearne correttamente il valore. Il
blocco viene rimosso chiamando \func{munlock}.

Altre due funzioni di sistema, \funcd{mlockall} e \funcd{munlockall},
consentono di bloccare genericamente la paginazione per l'intero spazio di
indirizzi di un processo.  I prototipi di queste funzioni sono:

\begin{funcproto}{ 
\fhead{sys/mman.h} 
\fdecl{int mlockall(int flags)}
\fdesc{Blocca la paginazione per lo spazio di indirizzi del processo corrente.} 
\fdecl{int munlockall(void)}
\fdesc{Sblocca la paginazione per lo spazio di indirizzi del processo corrente.}
}
{Codici di ritorno ed errori sono gli stessi di \func{mlock} e \func{munlock},
  tranne per \errcode{EINVAL} che viene restituito solo se si è specificato
  con \func{mlockall} un valore sconosciuto per \param{flags}.}
\end{funcproto}

L'argomento \param{flags} di \func{mlockall} permette di controllarne il
comportamento; esso deve essere specificato come maschera binaria dei valori
espressi dalle costanti riportate in tab.~\ref{tab:mlockall_flags}. 

\begin{table}[htb]
  \footnotesize
  \centering
  \begin{tabular}[c]{|l|p{8cm}|}
    \hline
    \textbf{Valore} & \textbf{Significato} \\
    \hline
    \hline
    \constd{MCL\_CURRENT}& blocca tutte le pagine correntemente mappate nello
                           spazio di indirizzi del processo.\\
    \constd{MCL\_FUTURE} & blocca tutte le pagine che verranno mappate nello
                           spazio di indirizzi del processo.\\
    \constd{MCL\_ONFAULT}& esegue il blocco delle pagine selezionate solo
                           quando vengono utilizzate (dal kernel 4.4).\\
   \hline
  \end{tabular}
  \caption{Valori e significato dell'argomento \param{flags} della funzione
    \func{mlockall}.}
  \label{tab:mlockall_flags}
\end{table}

Con \func{mlockall} si possono bloccare tutte le pagine mappate nello spazio
di indirizzi del processo, sia che comprendano il segmento di testo, di dati,
lo \textit{stack}, lo \textit{heap} e pure le funzioni di libreria chiamate, i
file mappati in memoria, i dati del kernel mappati in \textit{user space}, la
memoria condivisa.  L'uso dell'argomento \param{flags} permette di selezionare
con maggior finezza le pagine da bloccare, ad esempio usando
\const{MCL\_FUTURE} ci si può limitare a tutte le pagine allocate a partire
dalla chiamata della funzione, mentre \const{MCL\_CURRENT} blocca tutte quelle
correntemente mappate. L'uso di \func{munlockall} invece sblocca sempre tutte
le pagine di memoria correntemente mappate nello spazio di indirizzi del
programma.

A partire dal kernel 4.4 alla funzione \func{mlockall} è stato aggiunto un
altro flag, \const{MCL\_ONFAULT}, che può essere abbinato a entrambi gli altri
due flag, e consente di modificare il comportamento della funzione per
ottenere migliori prestazioni.

Il problema che si presenta infatti è che eseguire un \textit{memory lock} per
un intervallo ampio di memoria richiede che questa venga comunque allocata in
RAM, con altrettanti \textit{page fault} che ne assicurino la presenza; questo
vale per tutto l'intervallo e può avere un notevole costo in termini di
prestazioni, anche quando poi, nell'esecuzione del programma, venisse usata
solo una piccola parte dello stesso. L'uso di \const{MCL\_ONFAULT} previene il
\textit{page faulting} immediato di tutto l'intervallo, le pagine
dell'intervallo verranno bloccate, ma solo quando un \textit{page fault}
dovuto all'accesso ne richiede l'allocazione effettiva in RAM.

Questo stesso comportamento non è ottenibile con \func{mlock}, che non dispone
di un argomento \param{flag} che consenta di richiederlo, per questo sempre
con il kernel 4.4 è stata aggiunta una ulteriore funzione di sistema,
\funcd{mlock2}, il cui prototipo è:

\begin{funcproto}{
  \fhead{sys/mman.h} 
  \fdecl{int mlock2(const void *addr, size\_t len, int flags)}
  \fdesc{Blocca la paginazione su un intervallo di memoria.}
}
{Le funzione ritornano $0$ in caso di successo e $-1$ in caso di errore, nel
  qual caso \var{errno} assume gli stessi valori di \func{mlock} con
  l'aggiunta id un possibile \errcode{EINVAL} anche se si è indicato un valore
  errato di \param{flags}.}
\end{funcproto}

% NOTA: per mlock2, introdotta con il kernel 4.4 (vedi
% http://lwn.net/Articles/650538/)

Indicando un valore nullo per \param{flags} il comportamento della funzione è
identico a quello di \func{mlock}, l'unico altro valore possibile è
\constd{MLOCK\_ONFAULT} che ha lo stesso effetto sull'allocazione delle pagine
in RAM già descritto per \const{MCL\_ONFAULT}.

Si tenga presente che un processo \textit{real-time} che intende usare il
\textit{memory locking} con \func{mlockall} per prevenire l'avvenire di un
eventuale \textit{page fault} ed il conseguente rallentamento (probabilmente
inaccettabile) dei tempi di esecuzione, deve comunque avere delle accortezze.
In particolare si deve assicurare di aver preventivamente bloccato una
quantità di spazio nello \textit{stack} sufficiente a garantire l'esecuzione
di tutte le funzioni che hanno i requisiti di criticità sui tempi. Infatti,
anche usando \const{MCL\_FUTURE}, in caso di allocazione di una nuova pagina
nello \textit{stack} durante l'esecuzione di una funzione (precedentemente non
usata e quindi non bloccata) si potrebbe avere un \textit{page fault}.

In genere si ovvia a questa problematica chiamando inizialmente una funzione
che definisca una quantità sufficientemente ampia di variabili automatiche
(che si ricordi vengono allocate nello \textit{stack}) e ci scriva, in modo da
esser sicuri che le corrispondenti pagine vengano mappate nello spazio di
indirizzi del processo, per poi bloccarle. La scrittura è necessaria perché il
kernel usa il meccanismo di \textit{copy on write} (vedi
sez.~\ref{sec:proc_fork}) e le pagine potrebbero non essere allocate
immediatamente.

\itindend{memory~locking}
\index{memoria~virtuale|)} 


\subsection{Gestione avanzata dell'allocazione della memoria} 
\label{sec:proc_memory_adv_management}

La trattazione delle funzioni di allocazione di sez.~\ref{sec:proc_mem_alloc}
si è limitata a coprire le esigenze generiche di un programma, in cui non si
hanno dei requisiti specifici e si lascia il controllo delle modalità di
allocazione alle funzioni di libreria.  Tuttavia esistono una serie di casi in
cui può essere necessario avere un controllo più dettagliato delle modalità
con cui la memoria viene allocata; nel qual caso potranno venire in aiuto le
funzioni trattate in questa sezione.

Le prime funzioni che tratteremo sono quelle che consentono di richiedere di
allocare un blocco di memoria ``\textsl{allineato}'' ad un multiplo una certa
dimensione. Questo tipo di esigenza emerge usualmente quando si devono
allocare dei buffer da utilizzare per eseguire dell'I/O diretto su dispositivi
a blocchi. In questo caso infatti il trasferimento di dati viene eseguito per
blocchi di dimensione fissa, ed è richiesto che l'indirizzo di partenza del
buffer sia un multiplo intero di questa dimensione, usualmente 512 byte. In
tal caso l'uso di \func{malloc} non è sufficiente, ed occorre utilizzare una
funzione specifica.

Tradizionalmente per rispondere a questa esigenza sono state create due
funzioni diverse, \funcd{memalign} e \funcd{valloc}, oggi obsolete, cui si
aggiunge \funcd{pvalloc} come estensione GNU, anch'essa obsoleta; i rispettivi
prototipi sono:

\begin{funcproto}{ 
\fhead{malloc.h} 
\fdecl{void *valloc(size\_t size)}
\fdesc{Alloca un blocco di memoria allineato alla dimensione di una pagina di
  memoria.}  
\fdecl{void *memalign(size\_t boundary, size\_t size)}
\fdesc{Alloca un blocco di memoria allineato ad un multiplo
  di \param{boundary}.} 
\fdecl{void *pvalloc(size\_t size)}
\fdesc{Alloca un blocco di memoria allineato alla dimensione di una pagina di
  memoria.}  
}
{Entrambe le funzioni ritornano un puntatore al blocco di memoria allocato in
  caso di successo e \val{NULL} in caso di errore, nel qual caso \var{errno}
  assumerà uno dei valori:
  \begin{errlist}
  \item[\errcode{EINVAL}] \param{boundary} non è una potenza di due.
  \item[\errcode{ENOMEM}] non c'è memoria sufficiente per l'allocazione.
  \end{errlist}}
\end{funcproto}

Le funzioni restituiscono il puntatore al buffer di memoria allocata di
dimensioni pari a \param{size}, che per \func{memalign} sarà un multiplo di
\param{boundary} mentre per \func{valloc} un multiplo della dimensione di una
pagina di memoria; lo stesso vale per \func{pvalloc} che però arrotonda
automaticamente la dimensione dell'allocazione al primo multiplo di una
pagina. Nel caso della versione fornita dalla \acr{glibc} la memoria allocata
con queste funzioni deve essere liberata con \func{free}, cosa che non è detto
accada con altre implementazioni.

Nessuna delle due funzioni ha una chiara standardizzazione e nessuna delle due
compare in POSIX.1, inoltre ci sono indicazioni discordi sui file che ne
contengono la definizione;\footnote{secondo SUSv2 \func{valloc} è definita in
  \headfile{stdlib.h}, mentre sia la \acr{glibc} che le precedenti \acr{libc4}
  e \acr{libc5} la dichiarano in \headfile{malloc.h}, lo stesso vale per
  \func{memalign} che in alcuni sistemi è dichiarata in \headfile{stdlib.h}.}
per questo motivo il loro uso è sconsigliato, essendo state sostituite dalla
nuova \funcd{posix\_memalign}, che è stata standardizzata in POSIX.1d e
disponibile dalla \acr{glibc} 2.1.91; il suo prototipo è:

\begin{funcproto}{ 
\fhead{stdlib.h} 
\fdecl{posix\_memalign(void **memptr, size\_t alignment, size\_t size)}
\fdesc{Alloca un buffer di memoria allineato ad un multiplo
  di \param{alignment}.}   
}
{Entrambe le funzioni ritornano un puntatore al blocco di memoria allocato in
  caso di successo e \val{NULL} in caso di errore, nel qual caso \var{errno}
  assumerà uno dei valori:
  \begin{errlist}
  \item[\errcode{EINVAL}] \param{alignment} non è potenza di due o un multiplo
    di \code{sizeof(void *)}.
  \item[\errcode{ENOMEM}] non c'è memoria sufficiente per l'allocazione.
  \end{errlist}}
\end{funcproto}

La funzione restituisce il puntatore al buffer allocato di dimensioni pari
a \param{size} nella variabile (di tipo \texttt{void *}) posta all'indirizzo
indicato da \param{memptr}. La funzione fallisce nelle stesse condizioni delle
due funzioni precedenti, ma a loro differenza restituisce direttamente come
valore di ritorno il codice di errore.  Come per le precedenti la memoria
allocata con \func{posix\_memalign} deve essere disallocata con \func{free},
che in questo caso però è quanto richiesto dallo standard.

Dalla versione 2.16 della \acr{glibc} è stata aggiunta anche la funzione
\funcd{aligned\_alloc}, prevista dallo standard C11 (e disponibile definendo
\const{\_ISOC11\_SOURCE}), il cui prototipo è:

\begin{funcproto}{ 
\fhead{malloc.h} 
\fdecl{void *aligned\_alloc(size\_t alignment, size\_t size)}
\fdesc{Alloca un blocco di memoria allineato ad un multiplo
  di \param{alignment}.} 
}
{La funzione ha gli stessi valori di ritorno e codici di errore di
  \func{memalign}.}
\end{funcproto}

La funzione è identica a \func{memalign} ma richiede che \param{size} sia un
multiplo di \param{alignment}.  Infine si tenga presente infine che nessuna di
queste funzioni inizializza il buffer di memoria allocato, il loro
comportamento cioè è analogo, allineamento a parte, a quello di \func{malloc}.

Un secondo caso in cui risulta estremamente utile poter avere un maggior
controllo delle modalità di allocazione della memoria è quello in cui cercano
errori di programmazione. Esempi di questi errori sono i \textit{double free},
o i cosiddetti \itindex{buffer~overrun} \textit{buffer overrun}, cioè le
scritture su un buffer oltre le dimensioni della sua
allocazione,\footnote{entrambe queste operazioni causano in genere la
  corruzione dei dati di controllo delle funzioni di allocazione, che vengono
  anch'essi mantenuti nello \textit{heap} per tenere traccia delle zone di
  memoria allocata.} o i classici \textit{memory leak}.

Abbiamo visto in sez.~\ref{sec:proc_mem_lock} come una prima funzionalità di
ausilio nella ricerca di questi errori sia l'uso della variabile di ambiente
\envvar{MALLOC\_CHECK\_}.  Una modalità alternativa per effettuare dei
controlli di consistenza sullo stato delle allocazioni di memoria eseguite con
\func{malloc}, anche questa fornita come estensione specifica (e non standard)
della \acr{glibc}, è quella di utilizzare la funzione \funcd{mcheck}, che deve
essere chiamata prima di eseguire qualunque allocazione con \func{malloc}; il
suo prototipo è:

\begin{funcproto}{ 
\fhead{mcheck.h} 
\fdecl{int mcheck(void (*abortfn) (enum mcheck\_status status))}
\fdesc{Attiva i controlli di consistenza delle allocazioni di memoria.}   
}
{La funzione ritorna $0$ in caso di successo e $-1$ per un errore;
  \var{errno} non viene impostata.} 
\end{funcproto}

La funzione consente di registrare una funzione di emergenza che verrà
eseguita tutte le volte che, in una successiva esecuzione di \func{malloc},
venissero trovate delle inconsistenze, come delle operazioni di scrittura
oltre i limiti dei buffer allocati. Per questo motivo la funzione deve essere
chiamata prima di qualunque allocazione di memoria, altrimenti fallirà.

Se come primo argomento di \func{mcheck} si passa \val{NULL} verrà utilizzata
una funzione predefinita che stampa un messaggio di errore ed invoca la
funzione \func{abort} (vedi sez.~\ref{sec:sig_alarm_abort}), altrimenti si
dovrà creare una funzione personalizzata in grado di ricevere il tipo di
errore ed agire di conseguenza.

Nonostante la scarsa leggibilità del prototipo si tratta semplicemente di
definire una funzione di tipo \code{void abortfn(enum mcheck\_status status)},
che non deve restituire nulla e che deve avere un unico argomento di tipo
\code{mcheck\_status}. In caso di errore la funzione verrà eseguita ricevendo
un opportuno valore di \param{status} che è un tipo enumerato che può assumere
soltanto i valori di tab.~\ref{tab:mcheck_status_value} che indicano la
tipologia di errore riscontrata.

\begin{table}[htb]
  \centering
  \footnotesize
  \begin{tabular}[c]{|l|p{7cm}|}
    \hline
    \textbf{Valore} & \textbf{Significato} \\
    \hline
    \hline
    \constd{MCHECK\_OK}      & Riportato a \func{mprobe} se nessuna
                               inconsistenza è presente.\\
    \constd{MCHECK\_DISABLED}& Riportato a \func{mprobe} se si è chiamata
                               \func{mcheck} dopo aver già usato
                               \func{malloc}.\\
    \constd{MCHECK\_HEAD}    & I dati immediatamente precedenti il buffer sono
                               stati modificati, avviene in genere quando si
                               decrementa eccessivamente il valore di un
                               puntatore scrivendo poi prima dell'inizio del
                               buffer.\\
    \constd{MCHECK\_TAIL}    & I dati immediatamente seguenti il buffer sono
                               stati modificati, succede quando si va scrivere
                               oltre la dimensione corretta del buffer.\\
    \constd{MCHECK\_FREE}    & Il buffer è già stato disallocato.\\
    \hline
  \end{tabular}
  \caption{Valori dello stato dell'allocazione di memoria ottenibili dalla
    funzione di terminazione installata con \func{mcheck}.} 
  \label{tab:mcheck_status_value}
\end{table}

Una volta che si sia chiamata \func{mcheck} con successo si può anche
controllare esplicitamente lo stato delle allocazioni senza aspettare un
errore nelle relative funzioni utilizzando la funzione \funcd{mprobe}, il cui
prototipo è:

\begin{funcproto}{ 
\fhead{mcheck.h} 
\fdecl{enum mcheck\_status mprobe(ptr)}
\fdesc{Esegue un controllo di consistenza delle allocazioni.}   
}
{La funzione ritorna un codice fra quelli riportati in
   tab.~\ref{tab:mcheck_status_value} e non ha errori.} 
\end{funcproto}

La funzione richiede che si passi come argomento un puntatore ad un blocco di
memoria precedentemente allocato con \func{malloc} o \func{realloc}, e
restituisce lo stesso codice di errore che si avrebbe per la funzione di
emergenza ad una successiva chiamata di una funzione di allocazione, e poi i
primi due codici che indicano rispettivamente quando tutto è a posto o il
controllo non è possibile per non aver chiamato \func{mcheck} in tempo.

% TODO: trattare le altre funzionalità avanzate di \func{malloc}, mallopt,
% mtrace, muntrace, mallinfo e gli hook con le glibc 2.10 c'è pure malloc_info
% a sostituire mallinfo, vedi http://udrepper.livejournal.com/20948.html


\section{Argomenti, ambiente ed altre proprietà di un processo}
\label{sec:proc_options}

In questa sezione esamineremo le funzioni che permettono di gestire gli
argomenti e le opzioni, e quelle che consentono di manipolare ed utilizzare le
variabili di ambiente. Accenneremo infine alle modalità con cui si può gestire
la localizzazione di un programma modificandone il comportamento a seconda
della lingua o del paese a cui si vuole faccia riferimento nelle sue
operazioni. 

\subsection{Il formato degli argomenti}
\label{sec:proc_par_format}

Tutti i programmi hanno la possibilità di ricevere argomenti e opzioni quando
vengono lanciati e come accennato in sez.~\ref{sec:proc_main} questo viene
effettuato attraverso gli argomenti \param{argc} e \param{argv} ricevuti nella
funzione \code{main} all'avvio del programma. Questi argomenti vengono passati
al programma dalla shell o dal processo che esegue la \func{exec} (secondo le
modalità che vedremo in sez.~\ref{sec:proc_exec}) quando questo viene messo in
esecuzione.

Nel caso più comune il passaggio di argomenti ed opzioni viene effettuato
dalla shell, che si incarica di leggere la linea di comando con cui si lancia
il programma e di effettuarne la scansione (il cosiddetto \textit{parsing})
per individuare le parole che la compongono, ciascuna delle quali potrà essere
considerata un argomento o un'opzione. 

Di norma per individuare le parole che andranno a costituire la lista degli
argomenti viene usato come carattere di separazione lo spazio o il tabulatore,
ma la cosa dipende ovviamente dalle modalità con cui si effettua la scansione
e dalle convenzioni adottate dal programma che la esegue: ad esempio la shell
consente di proteggere con opportuni caratteri di controllo argomenti che
contengono degli spazi evitando di spezzarli in parole diverse.

\begin{figure}[htb]
  \centering
  \includegraphics[width=13cm]{img/argv_argc}
  % \begin{tikzpicture}[>=stealth]
  % \draw (0.5,2.5) rectangle (3.5,3);
  % \draw (2,2.75) node {\texttt{argc = 5}};
  % \draw (5,2.5) rectangle (8,3);
  % \draw (6.5,2.75) node {\texttt{argv[0]}};
  % \draw [->] (8,2.75) -- (9,2.75);
  % \draw (9,2.75) node [anchor=west] {\texttt{"touch"}};
  % \draw (5,2) rectangle (8,2.5);
  % \draw (6.5,2.25) node {\texttt{argv[1]}};
  % \draw [->] (8,2.25) -- (9,2.25);
  % \draw (9,2.25) node [anchor=west] {\texttt{"-r"}};
  % \draw (5,1.5) rectangle (8,2);
  % \draw (6.5,1.75) node {\texttt{argv[2]}};
  % \draw [->] (8,1.75) -- (9,1.75);
  % \draw (9,1.75) node [anchor=west] {\texttt{"riferimento.txt"}};
  % \draw (5,1.0) rectangle (8,1.5);
  % \draw (6.5,1.25) node {\texttt{argv[3]}};
  % \draw [->] (8,1.25) -- (9,1.25);
  % \draw (9,1.25) node [anchor=west] {\texttt{"-m"}};
  % \draw (5,0.5) rectangle (8,1.0);
  % \draw (6.5,0.75) node {\texttt{argv[4]}};
  % \draw [->] (8,0.75) -- (9,0.75);
  % \draw (9,0.75) node [anchor=west] {\texttt{"questofile.txt"}};
  % \draw (4.25,3.5) node{\texttt{"touch -r riferimento.txt -m questofile.txt"}};
  % \end{tikzpicture}
  \caption{Esempio dei valori di \param{argv} e \param{argc} generati nella 
    scansione di una riga di comando.}
  \label{fig:proc_argv_argc}
\end{figure}

Indipendentemente da come viene eseguita, il risultato finale della scansione
dovrà comunque essere la costruzione del vettore di puntatori \param{argv} in
cui si devono inserire in successione i puntatori alle stringhe costituenti i
vari argomenti ed opzioni da passare al programma, e della
variabile \param{argc} che deve essere inizializzata al numero di stringhe
contenute in \param{argv}. Nel caso della shell questo comporta ad esempio che
il primo argomento sia sempre il nome del programma. Un esempio di questo
meccanismo è mostrato in fig.~\ref{fig:proc_argv_argc}, che illustra il
risultato della scansione di una riga di comando.


\subsection{La gestione delle opzioni}
\label{sec:proc_opt_handling}

In generale un programma Unix riceve da linea di comando sia gli argomenti che
le opzioni, queste ultime sono standardizzate per essere riconosciute come
tali: un elemento di \param{argv} successivo al primo che inizia con il
carattere ``\texttt{-}'' e che non sia un singolo ``\texttt{-}'' o un
``\texttt{-{}-}'' viene considerato un'opzione.  In genere le opzioni sono
costituite da una lettera singola (preceduta dal carattere ``\texttt{-}'') e
possono avere o no un parametro associato.

Un esempio tipico può essere quello mostrato in
fig.~\ref{fig:proc_argv_argc}. In quel caso le opzioni sono \cmd{-r} e
\cmd{-m} e la prima vuole un parametro mentre la seconda no
(\cmd{questofile.txt} è un argomento del programma, non un parametro di
\cmd{-m}).

Per gestire le opzioni all'interno degli argomenti a linea di comando passati
in \param{argv} la libreria standard del C fornisce la funzione
\funcd{getopt}, che ha il seguente prototipo:

\begin{funcproto}{ 
\fhead{unistd.h} 
\fdecl{int getopt(int argc, char * const argv[], const char *optstring)}
\fdesc{Esegue la scansione delle opzioni negli argomenti della funzione
  \code{main}.} 
}
{Ritorna il carattere che segue l'opzione, ``\texttt{:}'' se manca un
  parametro all'opzione, ``\texttt{?}'' se l'opzione è sconosciuta, e $-1$ se
  non esistono altre opzioni.} 
\end{funcproto}

Questa funzione prende come argomenti le due variabili \param{argc} e
\param{argv} che devono essere quelle passate come argomenti di \code{main}
all'esecuzione del programma, ed una stringa \param{optstring} che indica
quali sono le opzioni valide. La funzione effettua la scansione della lista
degli argomenti ricercando ogni stringa che comincia con il carattere
``\texttt{-}'' e ritorna ogni volta che trova un'opzione valida.

La stringa \param{optstring} indica quali sono le opzioni riconosciute ed è
costituita da tutti i caratteri usati per identificare le singole opzioni, se
l'opzione ha un parametro al carattere deve essere fatto seguire il carattere
di due punti (``\texttt{:}''); nel caso di fig.~\ref{fig:proc_argv_argc} ad
esempio la stringa di opzioni avrebbe dovuto contenere \texttt{"r:m"}.

La modalità di uso di \func{getopt} è pertanto quella di chiamare più volte la
funzione all'interno di un ciclo, fintanto che essa non ritorna il valore $-1$
che indica che non ci sono più opzioni. Nel caso si incontri un'opzione non
dichiarata in \param{optstring} viene ritornato il carattere ``\texttt{?}''
mentre se un'opzione che lo richiede non è seguita da un parametro viene
ritornato il carattere ``\texttt{:}'', infine se viene incontrato il valore
``\texttt{-{}-}'' la scansione viene considerata conclusa, anche se vi sono
altri elementi di \param{argv} che cominciano con il carattere ``\texttt{-}''.

Quando \func{getopt} trova un'opzione fra quelle indicate in \param{optstring}
essa ritorna il valore numerico del carattere, in questo modo si possono
eseguire azioni specifiche usando uno \instruction{switch}; la funzione
inoltre inizializza alcune variabili globali:
\begin{itemize*}
\item \var{char *optarg} contiene il puntatore alla stringa parametro
  dell'opzione.
\item \var{int optind} alla fine della scansione restituisce l'indice del
  primo elemento di \param{argv} che non è un'opzione.
\item \var{int opterr} previene, se posto a zero, la stampa di un messaggio
  di errore in caso di riconoscimento di opzioni non definite.
\item \var{int optopt} contiene il carattere dell'opzione non riconosciuta.
\end{itemize*}

\begin{figure}[!htb]
  \footnotesize \centering
  \begin{minipage}[c]{\codesamplewidth}
  \includecodesample{listati/option_code.c}
  \end{minipage}
  \normalsize
  \caption{Esempio di codice per la gestione delle opzioni.}
  \label{fig:proc_options_code}
\end{figure}

In fig.~\ref{fig:proc_options_code} si è mostrata la sezione del programma
\file{fork\_test.c}, che useremo nel prossimo capitolo per effettuare dei test
sulla creazione dei processi, deputata alla decodifica delle opzioni a riga di
comando da esso supportate.

Si può notare che si è anzitutto (\texttt{\small 1}) disabilitata la stampa di
messaggi di errore per opzioni non riconosciute, per poi passare al ciclo per
la verifica delle opzioni (\texttt{\small 2-27}); per ciascuna delle opzioni
possibili si è poi provveduto ad un'azione opportuna, ad esempio per le tre
opzioni che prevedono un parametro si è effettuata la decodifica del medesimo,
il cui indirizzo è contenuto nella variabile \var{optarg}), avvalorando la
relativa variabile (\texttt{\small 12-14}, \texttt{\small 15-17} e
\texttt{\small 18-20}). Completato il ciclo troveremo in \var{optind}
l'indice in \code{argv[]} del primo degli argomenti rimanenti nella linea di
comando.

Normalmente \func{getopt} compie una permutazione degli elementi di
\param{argv} cosicché alla fine della scansione gli elementi che non sono
opzioni sono spostati in coda al vettore. Oltre a questa esistono altre due
modalità di gestire gli elementi di \param{argv}; se \param{optstring} inizia
con il carattere ``\texttt{+}'' (o è impostata la variabile di ambiente
\cmd{POSIXLY\_CORRECT}) la scansione viene fermata non appena si incontra un
elemento che non è un'opzione.

L'ultima modalità, usata quando un programma può gestire la mescolanza fra
opzioni e argomenti, ma se li aspetta in un ordine definito, si attiva
quando \param{optstring} inizia con il carattere ``\texttt{-}''. In questo caso
ogni elemento che non è un'opzione viene considerato comunque un'opzione e
associato ad un valore di ritorno pari ad 1, questo permette di identificare
gli elementi che non sono opzioni, ma non effettua il riordinamento del
vettore \param{argv}.


\subsection{Le variabili di ambiente}
\label{sec:proc_environ}

\index{variabili!di~ambiente|(}
Oltre agli argomenti passati a linea di comando esiste un'altra modalità che
permette di trasferire ad un processo delle informazioni in modo da
modificarne il comportamento.  Ogni processo infatti riceve dal sistema, oltre
alle variabili \param{argv} e \param{argc} anche un \textsl{ambiente} (in
inglese \textit{environment}); questo viene espresso nella forma di una lista
(chiamata \textit{environment list}) delle cosiddette \textsl{variabili di
  ambiente}, i valori di queste variabili possono essere poi usati dal
programma.

Anche in questo caso la lista delle \textsl{variabili di ambiente} deve essere
costruita ed utilizzata nella chiamata alla funzione \func{exec} (torneremo su
questo in sez.~\ref{sec:proc_exec}) quando questo viene lanciato. Come per la
lista degli argomenti anche questa lista è un vettore di puntatori a
caratteri, ciascuno dei quali punta ad una stringa, terminata da un
\val{NULL}. A differenza di \code{argv[]} in questo caso non si ha una
lunghezza del vettore data da un equivalente di \param{argc}, ma la lista è
terminata da un puntatore nullo.

L'indirizzo della lista delle variabili di ambiente è passato attraverso la
variabile globale \var{environ}, che viene definita automaticamente per
ciascun processo, e a cui si può accedere attraverso una semplice
dichiarazione del tipo:
\includecodesnip{listati/env_ptr.c}
un esempio della struttura di questa lista, contenente alcune delle variabili
più comuni che normalmente sono definite dal sistema, è riportato in
fig.~\ref{fig:proc_envirno_list}.
\begin{figure}[htb]
  \centering
  \includegraphics[width=13cm]{img/environ_var}
  % \begin{tikzpicture}[>=stealth]
  % \draw (2,3.5) node {\textsf{Environment pointer}};
  % \draw (6,3.5) node {\textsf{Environment list}};
  % \draw (10.5,3.5) node {\textsf{Environment string}};
  % \draw (0.5,2.5) rectangle (3.5,3);
  % \draw (2,2.75) node {\texttt{environ}};
  % \draw [->] (3.5,2.75) -- (4.5,2.75);
  % \draw (4.5,2.5) rectangle (7.5,3);
  % \draw (6,2.75) node {\texttt{environ[0]}};
  % \draw (4.5,2) rectangle (7.5,2.5);
  % \draw (6,2.25) node {\texttt{environ[1]}};
  % \draw (4.5,1.5) rectangle (7.5,2);
  % \draw (4.5,1) rectangle (7.5,1.5);
  % \draw (4.5,0.5) rectangle (7.5,1);
  % \draw (4.5,0) rectangle (7.5,0.5);
  % \draw (6,0.25) node {\texttt{NULL}};
  % \draw [->] (7.5,2.75) -- (8.5,2.75);
  % \draw (8.5,2.75) node[right] {\texttt{HOME=/home/piccardi}};
  % \draw [->] (7.5,2.25) -- (8.5,2.25);
  % \draw (8.5,2.25) node[right] {\texttt{PATH=:/bin:/usr/bin}};
  % \draw [->] (7.5,1.75) -- (8.5,1.75);
  % \draw (8.5,1.75) node[right] {\texttt{SHELL=/bin/bash}};
  % \draw [->] (7.5,1.25) -- (8.5,1.25);
  % \draw (8.5,1.25) node[right] {\texttt{EDITOR=emacs}};
  % \draw [->] (7.5,0.75) -- (8.5,0.75);
  % \draw (8.5,0.75) node[right] {\texttt{OSTYPE=linux-gnu}};
  % \end{tikzpicture}
  \caption{Esempio di lista delle variabili di ambiente.}
  \label{fig:proc_envirno_list}
\end{figure}

Per convenzione le stringhe che definiscono l'ambiente sono tutte del tipo
\textsl{\texttt{NOME=valore}} ed in questa forma che le funzioni di gestione
che vedremo a breve se le aspettano, se pertanto si dovesse costruire
manualmente un ambiente si abbia cura di rispettare questa convenzione.
Inoltre alcune variabili, come quelle elencate in
fig.~\ref{fig:proc_envirno_list}, sono definite dal sistema per essere usate
da diversi programmi e funzioni: per queste c'è l'ulteriore convenzione di
usare nomi espressi in caratteri maiuscoli.\footnote{ma si tratta solo di una
  convenzione, niente vieta di usare caratteri minuscoli, come avviene in vari
  casi.}

Il kernel non usa mai queste variabili, il loro uso e la loro interpretazione è
riservata alle applicazioni e ad alcune funzioni di libreria; in genere esse
costituiscono un modo comodo per definire un comportamento specifico senza
dover ricorrere all'uso di opzioni a linea di comando o di file di
configurazione. É di norma cura della shell, quando esegue un comando, passare
queste variabili al programma messo in esecuzione attraverso un uso opportuno
delle relative chiamate (si veda sez.~\ref{sec:proc_exec}).

La shell ad esempio ne usa molte per il suo funzionamento, come \envvar{PATH}
per indicare la lista delle directory in cui effettuare la ricerca dei comandi
o \envvar{PS1} per impostare il proprio \textit{prompt}. Alcune di esse, come
\envvar{HOME}, \envvar{USER}, ecc. sono invece definite al login (per i
dettagli si veda sez.~\ref{sec:sess_login}), ed in genere è cura della propria
distribuzione definire le opportune variabili di ambiente in uno script di
avvio. Alcune servono poi come riferimento generico per molti programmi, come
\envvar{EDITOR} che indica l'editor preferito da invocare in caso di
necessità. Una in particolare, \envvar{LANG}, serve a controllare la
localizzazione del programma 
%(su cui torneremo in sez.~\ref{sec:proc_localization}) 
per adattarlo alla lingua ed alle convezioni
dei vari paesi.

Gli standard POSIX e XPG3 definiscono alcune di queste variabili (le più
comuni), come riportato in tab.~\ref{tab:proc_env_var}. GNU/Linux le supporta
tutte e ne definisce anche altre, in particolare poi alcune funzioni di
libreria prevedono la presenza di specifiche variabili di ambiente che ne
modificano il comportamento, come quelle usate per indicare una localizzazione
e quelle per indicare un fuso orario; una lista più completa che comprende
queste ed ulteriori variabili si può ottenere con il comando \cmd{man 7
  environ}.

\begin{table}[htb]
  \centering
  \footnotesize
  \begin{tabular}[c]{|l|c|c|c|l|}
    \hline
    \textbf{Variabile} & \textbf{POSIX} & \textbf{XPG3} 
    & \textbf{Linux} & \textbf{Descrizione} \\
    \hline
    \hline
    \texttt{USER}   &$\bullet$&$\bullet$&$\bullet$& Nome utente.\\
    \texttt{LOGNAME}&$\bullet$&$\bullet$&$\bullet$& Nome di login.\\
    \texttt{HOME}   &$\bullet$&$\bullet$&$\bullet$& Directory base
                                                    dell'utente.\\
    \texttt{LANG}   &$\bullet$&$\bullet$&$\bullet$& Localizzazione.\\
    \texttt{PATH}   &$\bullet$&$\bullet$&$\bullet$& Elenco delle directory
                                                    dei programmi.\\
    \texttt{PWD}    &$\bullet$&$\bullet$&$\bullet$& Directory corrente.\\
    \texttt{SHELL}  &$\bullet$&$\bullet$&$\bullet$& Shell in uso.\\
    \texttt{TERM}   &$\bullet$&$\bullet$&$\bullet$& Tipo di terminale.\\
    \texttt{PAGER}  &$\bullet$&$\bullet$&$\bullet$& Programma per vedere i
                                                    testi.\\
    \texttt{EDITOR} &$\bullet$&$\bullet$&$\bullet$& Editor preferito.\\
    \texttt{BROWSER}&$\bullet$&$\bullet$&$\bullet$& Browser preferito.\\
    \texttt{TMPDIR} &$\bullet$&$\bullet$&$\bullet$& Directory dei file
                                                    temporanei.\\
    \hline
  \end{tabular}
  \caption{Esempi delle variabili di ambiente più comuni definite da vari
    standard.} 
  \label{tab:proc_env_var}
\end{table}

Lo standard ANSI C prevede l'esistenza di un ambiente, e pur non entrando
nelle specifiche di come sono strutturati i contenuti, definisce la funzione
\funcd{getenv} che permette di ottenere i valori delle variabili di ambiente;
il suo prototipo è:

\begin{funcproto}{ 
\fhead{stdlib.h}
\fdecl{char *getenv(const char *name)}
\fdesc{Cerca una variabile di ambiente del processo.} 
}
{La funzione ritorna il puntatore alla stringa contenente il valore della
  variabile di ambiente in caso di successo e \val{NULL} per un errore.} 
\end{funcproto}

La funzione effettua una ricerca nell'ambiente del processo cercando una
variabile il cui nome corrisponda a quanto indicato con
l'argomento \param{name}, ed in caso di successo ritorna il puntatore alla
stringa che ne contiene il valore, nella forma ``\texttt{NOME=valore}''.

\begin{table}[htb]
  \centering
  \footnotesize
  \begin{tabular}[c]{|l|c|c|c|c|c|c|}
    \hline
    \textbf{Funzione} & \textbf{ANSI C} & \textbf{POSIX.1} & \textbf{XPG3} & 
    \textbf{SVr4} & \textbf{BSD} & \textbf{Linux} \\
    \hline
    \hline
    \func{getenv}  & $\bullet$ & $\bullet$ & $\bullet$ 
                   & $\bullet$ & $\bullet$ & $\bullet$ \\
    \func{setenv}  &    --     &    --     &   --      
                   &    --     & $\bullet$ & $\bullet$ \\
    \func{unsetenv}&    --     &    --     &   --       
                   &    --     & $\bullet$ & $\bullet$ \\
    \func{putenv}  &    --     & opz.      & $\bullet$ 
                   &    --     & $\bullet$ & $\bullet$ \\
    \func{clearenv}&    --     & opz.      &   --
                   &    --     &    --     & $\bullet$ \\
    \hline
  \end{tabular}
  \caption{Funzioni per la gestione delle variabili di ambiente.}
  \label{tab:proc_env_func}
\end{table}

Oltre a questa funzione di lettura, che è l'unica definita dallo standard ANSI
C, nell'evoluzione dei sistemi Unix ne sono state proposte altre, da
utilizzare per impostare, modificare e cancellare le variabili di
ambiente. Uno schema delle funzioni previste nei vari standard e disponibili
in Linux è riportato in tab.~\ref{tab:proc_env_func}. Tutte le funzioni sono
state comunque inserite nello standard POSIX.1-2001, ad eccetto di
\func{clearenv} che è stata rigettata.

In Linux sono definite tutte le funzioni elencate in
tab.~\ref{tab:proc_env_func},\footnote{in realtà nelle libc4 e libc5 sono
  definite solo le prime quattro, \func{clearenv} è stata introdotta con la
  \acr{glibc} 2.0.} anche se parte delle funzionalità sono ridondanti. La
prima funzione di manipolazione che prenderemo in considerazione è
\funcd{putenv}, che consente di aggiungere, modificare e cancellare una
variabile di ambiente; il suo prototipo è:

\begin{funcproto}{ 
\fdecl{int putenv(char *string)}
\fdesc{Inserisce, modifica o rimuove una variabile d'ambiente.} 
}
{La funzione ritorna $0$ in caso di successo e $-1$ per un errore, che può
  essere solo \errval{ENOMEM}.}
\end{funcproto}

La funzione prende come argomento una stringa analoga a quella restituita da
\func{getenv} e sempre nella forma ``\texttt{NOME=valore}''. Se la variabile
specificata (nel caso \texttt{NOME}) non esiste la stringa sarà aggiunta
all'ambiente, se invece esiste il suo valore sarà impostato a quello
specificato dal contenuto di \param{string} (nel caso \texttt{valore}).  Se
invece si passa come argomento solo il nome di una variabile di ambiente
(cioè \param{string} è nella forma ``\texttt{NOME}'' e non contiene il
carattere ``\texttt{=}'') allora questa, se presente nell'ambiente, verrà
cancellata.

Si tenga presente che, seguendo lo standard SUSv2, le \acr{glibc} successive
alla versione 2.1.2 aggiungono direttamente \param{string} nella lista delle
variabili di ambiente illustrata in fig.~\ref{fig:proc_envirno_list}
sostituendo il relativo puntatore;\footnote{il comportamento è lo stesso delle
  vecchie \acr{libc4} e \acr{libc5}; nella \acr{glibc}, dalla versione 2.0
  alla 2.1.1, veniva invece fatta una copia, seguendo il comportamento di
  BSD4.4; dato che questo può dar luogo a perdite di memoria e non rispetta lo
  standard il comportamento è stato modificato a partire dalla 2.1.2,
  eliminando anche, sempre in conformità a SUSv2, l'attributo \direct{const}
  dal prototipo.}  pertanto ogni cambiamento alla stringa in questione si
riflette automaticamente sull'ambiente, e quindi si deve evitare di passare a
questa funzione una variabile automatica (per evitare i problemi esposti in
sez.~\ref{sec:proc_var_passing}). Benché non sia richiesto dallo standard,
nelle versioni della \acr{glibc} a partire dalla 2.1 la funzione è rientrante
(vedi sez.~\ref{sec:proc_reentrant}).

Infine quando una chiamata a \func{putenv} comporta la necessità di creare una
nuova versione del vettore \var{environ} questo sarà allocato automaticamente,
ma la versione corrente sarà deallocata solo se anch'essa è risultante da
un'allocazione fatta in precedenza da un'altra \func{putenv}. Questo avviene
perché il vettore delle variabili di ambiente iniziale, creato dalla chiamata
ad \func{exec} (vedi sez.~\ref{sec:proc_exec}) è piazzato nella memoria al di
sopra dello \textit{stack}, (vedi fig.~\ref{fig:proc_mem_layout}) e non nello
\textit{heap} e quindi non può essere deallocato.  Inoltre la memoria
associata alle variabili di ambiente eliminate non viene liberata.

Come alternativa a \func{putenv} si può usare la funzione \funcd{setenv} che
però consente solo di aggiungere o modificare una variabile di ambiente; il
suo prototipo è:

\begin{funcproto}{ 
\fhead{stdlib.h}
\fdecl{int setenv(const char *name, const char *value, int overwrite)}
\fdesc{Inserisce o modifica una variabile di ambiente.} 
}
{La funzione ritorna $0$ in caso di successo e $-1$ per un errore,
  nel qual caso \var{errno} assumerà uno dei valori:
  \begin{errlist}
  \item[\errcode{EINVAL}] \param{name} è \val{NULL} o una stringa di lunghezza
  nulla o che contiene il carattere ``\texttt{=}''.
  \item[\errcode{ENOMEM}] non c'è memoria sufficiente per aggiungere una nuova
    variabile all'ambiente.
\end{errlist}}
\end{funcproto}

La funzione consente di specificare separatamente nome e valore della
variabile di ambiente da aggiungere negli argomenti \param{name}
e \param{value}. Se la variabile è già presente nell'ambiente
l'argomento \param{overwrite} specifica il comportamento della funzione, se
diverso da zero sarà sovrascritta, se uguale a zero sarà lasciata immutata.  A
differenza di \func{putenv} la funzione esegue delle copie del contenuto degli
argomenti \param{name} e \param{value} e non è necessario preoccuparsi di
allocarli in maniera permanente.

La cancellazione di una variabile di ambiente viene invece gestita
esplicitamente con \funcd{unsetenv}, il cui prototipo è:

\begin{funcproto}{ 
\fhead{stdlib.h}
\fdecl{int unsetenv(const char *name)}
\fdesc{Rimuove una variabile di ambiente.} 
}
{La funzione ritorna $0$ in caso di successo e $-1$ per un errore,
  nel qual caso \var{errno} assumerà uno dei valori:
  \begin{errlist}
  \item[\errcode{EINVAL}] \param{name} è \val{NULL} o una stringa di lunghezza
  nulla o che contiene il carattere ``\texttt{=}''.
\end{errlist}}
\end{funcproto}

La funzione richiede soltanto il nome della variabile di ambiente
nell'argomento \param{name}, se la variabile non esiste la funzione ritorna
comunque con un valore di successo.\footnote{questo con le versioni della
  \acr{glibc} successive la 2.2.2, per le precedenti \func{unsetenv} era
  definita come \texttt{void} e non restituiva nessuna informazione.}

L'ultima funzione per la gestione dell'ambiente è
\funcd{clearenv},\footnote{che come accennato è l'unica non presente nello
  standard POSIX.1-2000, ed è disponibili solo per versioni della \acr{glibc}
  a partire dalla 2.0; per poterla utilizzare occorre aver definito le macro
  \macro{\_SVID\_SOURCE} e \macro{\_XOPEN\_SOURCE}.} che viene usata per
cancellare completamente tutto l'ambiente; il suo prototipo è:

\begin{funcproto}{ 
\fhead{stdlib.h}
\fdecl{int clearenv(void)}
\fdesc{Cancella tutto l'ambiente.} 
}
{La funzione ritorna $0$ in caso di successo e un valore diverso da zero per
  un errore.}
\end{funcproto}

In genere si usa questa funzione in maniera precauzionale per evitare i
problemi di sicurezza connessi nel trasmettere ai programmi che si invocano un
ambiente che può contenere dei dati non controllati, le cui variabili possono
causare effetti indesiderati. Con l'uso della funzione si provvede alla
cancellazione di tutto l'ambiente originale in modo da poterne costruirne una
versione ``\textsl{sicura}'' da zero.

\index{variabili!di~ambiente|)}


% \subsection{La localizzazione}
% \label{sec:proc_localization}

% Abbiamo accennato in sez.~\ref{sec:proc_environ} come la variabile di ambiente
% \envvar{LANG} sia usata per indicare ai processi il valore della cosiddetta
% \textsl{localizzazione}. Si tratta di una funzionalità fornita dalle librerie
% di sistema\footnote{prenderemo in esame soltanto il caso della \acr{glibc}.}
% che consente di gestire in maniera automatica sia la lingua in cui vengono
% stampati i vari messaggi (come i messaggi associati agli errori che vedremo in
% sez.~\ref{sec:sys_strerror}) che le convenzioni usate nei vari paesi per una
% serie di aspetti come il formato dell'ora, quello delle date, gli ordinamenti
% alfabetici, le espressioni della valute, ecc.

% Da finire.

% La localizzazione di un programma si può selezionare con la 

% In realtà perché un programma sia effettivamente localizzato non è sufficiente 

% TODO trattare, quando ci sarà tempo, setlocale ed il resto


%\subsection{Opzioni in formato esteso}
%\label{sec:proc_opt_extended}

%Oltre alla modalità ordinaria di gestione delle opzioni trattata in
%sez.~\ref{sec:proc_opt_handling} la \acr{glibc} fornisce una modalità
%alternativa costituita dalle cosiddette \textit{long-options}, che consente di
%esprimere le opzioni in una forma più descrittiva che nel caso più generale è
%qualcosa del tipo di ``\texttt{-{}-option-name=parameter}''.

%(NdA: questa parte verrà inserita in seguito).

% TODO opzioni in formato esteso

% TODO trattare il vettore ausiliario e getauxval (vedi
% http://lwn.net/Articles/519085/)


\section{Problematiche di programmazione generica}
\label{sec:proc_gen_prog}

Benché questo non sia un libro sul linguaggio C, è opportuno affrontare alcune
delle problematiche generali che possono emergere nella programmazione con
questo linguaggio e di quali precauzioni o accorgimenti occorre prendere per
risolverle. Queste problematiche non sono specifiche di sistemi unix-like o
\textit{multitasking}, ma avendo trattato in questo capitolo il comportamento
dei processi visti come entità a sé stanti, le riportiamo qui.


\subsection{Il passaggio di variabili e valori di ritorno nelle funzioni}
\label{sec:proc_var_passing}

Una delle caratteristiche standard del C è che le variabili vengono passate
alle funzioni che si invocano in un programma attraverso un meccanismo che
viene chiamato \textit{by value}, diverso ad esempio da quanto avviene con il
Fortran, dove le variabili sono passate, come suol dirsi, \textit{by
  reference}, o dal C++ dove la modalità del passaggio può essere controllata
con l'operatore \cmd{\&}.

Il passaggio di una variabile \textit{by value} significa che in realtà quello
che viene passato alla funzione è una copia del valore attuale di quella
variabile, copia che la funzione potrà modificare a piacere, senza che il
valore originale nella funzione chiamante venga toccato. In questo modo non
occorre preoccuparsi di eventuali effetti delle operazioni svolte nella
funzione stessa sulla variabile passata come argomento.

Questo però va inteso nella maniera corretta. Il passaggio \textit{by value}
vale per qualunque variabile, puntatori compresi; quando però in una funzione
si usano dei puntatori (ad esempio per scrivere in un buffer) in realtà si va
a modificare la zona di memoria a cui essi puntano, per cui anche se i
puntatori sono copie, i dati a cui essi puntano saranno sempre gli stessi, e
le eventuali modifiche avranno effetto e saranno visibili anche nella funzione
chiamante.

Nella maggior parte delle funzioni di libreria e delle \textit{system call} i
puntatori vengono usati per scambiare dati (attraverso i buffer o le strutture
a cui fanno riferimento) e le variabili normali vengono usate per specificare
argomenti; in genere le informazioni a riguardo dei risultati vengono passate
alla funzione chiamante attraverso il valore di ritorno.  È buona norma
seguire questa pratica anche nella programmazione normale.

\itindbeg{value~result~argument}

Talvolta però è necessario che la funzione possa restituire indietro alla
funzione chiamante un valore relativo ad uno dei suoi argomenti usato anche in
ingresso.  Per far questo si usa il cosiddetto \textit{value result argument},
si passa cioè, invece di una normale variabile, un puntatore alla stessa. Gli
esempi di questa modalità di passaggio sono moltissimi, ad esempio essa viene
usata nelle funzioni che gestiscono i socket (in
sez.~\ref{sec:TCP_functions}), in cui, per permettere al kernel di restituire
informazioni sulle dimensioni delle strutture degli indirizzi utilizzate,
viene usato proprio questo meccanismo.

Occorre tenere ben presente questa differenza, perché le variabili passate in
maniera ordinaria, che vengono inserite nello \textit{stack}, cessano di
esistere al ritorno di una funzione, ed ogni loro eventuale modifica
all'interno della stessa sparisce con la conclusione della stessa, per poter
passare delle informazioni occorre quindi usare un puntatore che faccia
riferimento ad un indirizzo accessibile alla funzione chiamante.

\itindend{value~result~argument}

Questo requisito di accessibilità è fondamentale, infatti dei possibili
problemi che si possono avere con il passaggio dei dati è quello di restituire
alla funzione chiamante dei dati che sono contenuti in una variabile
automatica.  Ovviamente quando la funzione ritorna la sezione dello
\textit{stack} che conteneva la variabile automatica (si ricordi quanto detto
in sez.~\ref{sec:proc_mem_alloc}) verrà liberata automaticamente e potrà
essere riutilizzata all'invocazione di un'altra funzione, con le immaginabili
conseguenze, quasi invariabilmente catastrofiche, di sovrapposizione e
sovrascrittura dei dati.

Per questo una delle regole fondamentali della programmazione in C è che
all'uscita di una funzione non deve restare nessun riferimento alle sue
variabili locali. Qualora sia necessario utilizzare delle variabili che devono
essere viste anche dalla funzione chiamante queste devono essere allocate
esplicitamente, o in maniera statica usando variabili globali o dichiarate
come \direct{extern},\footnote{la direttiva \direct{extern} informa il
  compilatore che la variabile che si è dichiarata in una funzione non è da
  considerarsi locale, ma globale, e per questo allocata staticamente e
  visibile da tutte le funzioni dello stesso programma.} o dinamicamente con
una delle funzioni della famiglia \func{malloc}, passando opportunamente il
relativo puntatore fra le funzioni.


\subsection{Il passaggio di un numero variabile di argomenti}
\label{sec:proc_variadic}

\index{funzioni!\textit{variadic}|(}

Come vedremo nei capitoli successivi, non sempre è possibile specificare un
numero fisso di argomenti per una funzione.  Lo standard ISO C prevede nella
sua sintassi la possibilità di definire delle \textit{variadic function} che
abbiano un numero variabile di argomenti, attraverso l'uso nella dichiarazione
della funzione dello speciale costrutto ``\texttt{...}'', che viene chiamato
\textit{ellipsis}.

Lo standard però non provvede a livello di linguaggio alcun meccanismo con cui
dette funzioni possono accedere ai loro argomenti.  L'accesso viene pertanto
realizzato a livello della libreria standard del C che fornisce gli strumenti
adeguati.  L'uso di una \textit{variadic function} prevede quindi tre punti:
\begin{itemize*}
\item \textsl{dichiarare} la funzione come \textit{variadic} usando un
  prototipo che contenga una \textit{ellipsis};
\item \textsl{definire} la funzione come \textit{variadic} usando la stessa
  \textit{ellipsis}, ed utilizzare le apposite macro che consentono la
  gestione di un numero variabile di argomenti;
\item \textsl{invocare} la funzione specificando prima gli argomenti fissi, ed
  a seguire quelli addizionali.
\end{itemize*}

Lo standard ISO C prevede che una \textit{variadic function} abbia sempre
almeno un argomento fisso. Prima di effettuare la dichiarazione deve essere
incluso l'apposito \textit{header file} \headfile{stdarg.h}; un esempio di
dichiarazione è il prototipo della funzione \func{execl} che vedremo in
sez.~\ref{sec:proc_exec}:
\includecodesnip{listati/exec_sample.c}
in questo caso la funzione prende due argomenti fissi ed un numero variabile
di altri argomenti, che andranno a costituire gli elementi successivi al primo
del vettore \param{argv} passato al nuovo processo. Lo standard ISO C richiede
inoltre che l'ultimo degli argomenti fissi sia di tipo
\textit{self-promoting}\footnote{il linguaggio C prevede che quando si
  mescolano vari tipi di dati, alcuni di essi possano essere \textsl{promossi}
  per compatibilità; ad esempio i tipi \ctyp{float} vengono convertiti
  automaticamente a \ctyp{double} ed i \ctyp{char} e gli \ctyp{short} ad
  \ctyp{int}. Un tipo \textit{self-promoting} è un tipo che verrebbe promosso
  a sé stesso.} il che esclude vettori, puntatori a funzioni e interi di tipo
\ctyp{char} o \ctyp{short} (con segno o meno). Una restrizione ulteriore di
alcuni compilatori è di non dichiarare l'ultimo argomento fisso come variabile
di tipo \direct{register}.\footnote{la direttiva \direct{register} del
  compilatore chiede che la variabile dichiarata tale sia mantenuta, nei
  limiti del possibile, all'interno di un registro del processore; questa
  direttiva è originaria dell'epoca dai primi compilatori, quando stava al
  programmatore scrivere codice ottimizzato, riservando esplicitamente alle
  variabili più usate l'uso dei registri del processore, oggi questa direttiva
  è in disuso pressoché completo dato che tutti i compilatori sono normalmente
  in grado di valutare con maggior efficacia degli stessi programmatori quando
  sia il caso di eseguire questa ottimizzazione.}

Una volta dichiarata la funzione il secondo passo è accedere ai vari argomenti
quando la si va a definire. Gli argomenti fissi infatti hanno un loro nome, ma
quelli variabili vengono indicati in maniera generica dalla
\textit{ellipsis}. L'unica modalità in cui essi possono essere recuperati è
pertanto quella sequenziale, in cui vengono estratti dallo \textit{stack}
secondo l'ordine in cui sono stati scritti nel prototipo della funzione.

\macrobeg{va\_start}

Per fare questo in \headfile{stdarg.h} sono definite delle macro specifiche,
previste dallo standard ISO C89, che consentono di eseguire questa operazione.
La prima di queste macro è \macro{va\_start}, che inizializza opportunamente
una lista degli argomenti, la sua definizione è:

{\centering
\begin{funcbox}{ 
\fhead{stdarg.h}
\fdecl{void va\_start(va\_list ap, last)}
\fdesc{Inizializza una lista degli argomenti di una funzione
  \textit{variadic}.} 
}
\end{funcbox}}

La macro inizializza il puntatore alla lista di argomenti \param{ap} che deve
essere una apposita variabile di tipo \type{va\_list}; il
parametro \param{last} deve indicare il nome dell'ultimo degli argomenti fissi
dichiarati nel prototipo della funzione \textit{variadic}.

\macrobeg{va\_arg}

La seconda macro di gestione delle liste di argomenti di una funzione
\textit{variadic} è \macro{va\_arg}, che restituisce in successione un
argomento della lista; la sua definizione è:

{\centering
\begin{funcbox}{ 
\fhead{stdarg.h}
\fdecl{type va\_arg(va\_list ap, type)}
\fdesc{Restituisce il valore del successivo argomento opzionale.} 
}
\end{funcbox}}
 
La macro restituisce il valore di un argomento, modificando opportunamente la
lista \param{ap} perché una chiamata successiva restituisca l'argomento
seguente. La macro richiede che si specifichi il tipo dell'argomento che si
andrà ad estrarre attraverso il parametro \param{type} che sarà anche il tipo
del valore da essa restituito. Si ricordi che il tipo deve essere
\textit{self-promoting}.

In generale è perfettamente legittimo richiedere meno argomenti di quelli che
potrebbero essere stati effettivamente forniti, per cui nella esecuzione delle
\macro{va\_arg} ci si può fermare in qualunque momento ed i restanti argomenti
saranno ignorati. Se invece si richiedono più argomenti di quelli
effettivamente forniti si otterranno dei valori indefiniti. Si avranno
risultati indefiniti anche quando si chiama \macro{va\_arg} specificando un
tipo che non corrisponde a quello usato per il corrispondente argomento.

\macrobeg{va\_end}

Infine una volta completata l'estrazione occorre indicare che si sono concluse
le operazioni con la macro \macrod{va\_end}, la cui definizione è:

{\centering
\begin{funcbox}{ 
\fhead{stdarg.h}
\fdecl{void va\_end(va\_list ap)}
\fdesc{Conclude l'estrazione degli argomenti di una funzione
  \textit{variadic}.} 
}
\end{funcbox}}
 
Dopo l'uso di \macro{va\_end} la variabile \param{ap} diventa indefinita e
successive chiamate a \macro{va\_arg} non funzioneranno.  Nel caso del
\cmd{gcc} l'uso di \macro{va\_end} può risultare inutile, ma è comunque
necessario usarla per chiarezza del codice, per compatibilità con diverse
implementazioni e per eventuali modifiche future a questo comportamento.

Riassumendo la procedura da seguire per effettuare l'estrazione degli
argomenti di una funzione \textit{variadic} è la seguente:
\begin{enumerate*}
\item inizializzare una lista degli argomenti attraverso la macro
  \macro{va\_start};
\item accedere agli argomenti con chiamate successive alla macro
  \macro{va\_arg}: la prima chiamata restituirà il primo argomento, la seconda
  il secondo e così via;
\item dichiarare la conclusione dell'estrazione degli argomenti invocando la
  macro \macro{va\_end}.
\end{enumerate*}

Si tenga presente che si possono usare anche più liste degli argomenti,
ciascuna di esse andrà inizializzata con \macro{va\_start} e letta con
\macro{va\_arg}, e ciascuna potrà essere usata per scandire la lista degli
argomenti in modo indipendente. Infine ciascuna scansione dovrà essere
terminata con \macro{va\_end}.

Un limite di queste macro è che i passi 1) e 3) devono essere eseguiti nel
corpo principale della funzione, il passo 2) invece può essere eseguito anche
in un'altra funzione, passandole lista degli argomenti \param{ap}. In questo
caso però al ritorno della funzione \macro{va\_arg} non può più essere usata
(anche se non si era completata l'estrazione) dato che il valore di \param{ap}
risulterebbe indefinito.

\macroend{va\_start}
\macroend{va\_arg}
\macroend{va\_end}

Esistono dei casi in cui è necessario eseguire più volte la scansione degli
argomenti e poter memorizzare una posizione durante la stessa. In questo caso
sembrerebbe naturale copiarsi la lista degli argomenti \param{ap} con una
semplice assegnazione ad un'altra variabile dello stesso tipo. Dato che una
delle realizzazioni più comuni di \type{va\_list} è quella di un puntatore
nello \textit{stack} all'indirizzo dove sono stati salvati gli argomenti, è
assolutamente normale pensare di poter effettuare questa operazione.

\index{tipo!opaco|(}

In generale però possono esistere anche realizzazioni diverse, ed è per questo
motivo che invece che un semplice puntatore, \typed{va\_list} è quello che
viene chiamato un \textsl{tipo opaco}. Si chiamano così quei tipi di dati, in
genere usati da una libreria, la cui struttura interna non deve essere vista
dal programma chiamante (da cui deriva il nome opaco) che li devono utilizzare
solo attraverso dalle opportune funzioni di gestione.

\index{tipo!opaco|)}

Per questo motivo una variabile di tipo \typed{va\_list} non può essere
assegnata direttamente ad un'altra variabile dello stesso tipo, ma lo standard
ISO C99\footnote{alcuni sistemi che non hanno questa macro provvedono al suo
  posto \macrod{\_\_va\_copy} che era il nome proposto in una bozza dello
  standard.}  ha previsto una macro ulteriore che permette di eseguire la
copia di una lista degli argomenti:

{\centering
\begin{funcbox}{ 
\fhead{stdarg.h}
\fdecl{void va\_copy(va\_list dest, va\_list src)}
\fdesc{Copia la lista degli argomenti di una funzione \textit{variadic}.} 
}
\end{funcbox}}

La macro copia l'attuale della lista degli argomenti \param{src} su una nuova
lista \param{dest}. Anche in questo caso è buona norma chiudere ogni
esecuzione di una \macrod{va\_copy} con una corrispondente \macro{va\_end} sul
nuovo puntatore alla lista degli argomenti.

La chiamata di una funzione con un numero variabile di argomenti, posto che la
si sia dichiarata e definita come tale, non prevede nulla di particolare;
l'invocazione è identica alle altre, con gli argomenti, sia quelli fissi che
quelli opzionali, separati da virgole. Quello che però è necessario tenere
presente è come verranno convertiti gli argomenti variabili.

In Linux gli argomenti dello stesso tipo sono passati allo stesso modo, sia
che siano fissi sia che siano opzionali (alcuni sistemi trattano diversamente
gli opzionali), ma dato che il prototipo non può specificare il tipo degli
argomenti opzionali, questi verranno sempre promossi, pertanto nella ricezione
dei medesimi occorrerà tenerne conto (ad esempio un \ctyp{char} verrà visto da
\macro{va\_arg} come \ctyp{int}).

Un altro dei problemi che si devono affrontare con le funzioni con un numero
variabile di argomenti è che non esiste un modo generico che permetta di
stabilire quanti sono gli argomenti effettivamente passati in una chiamata.

Esistono varie modalità per affrontare questo problema; una delle più
immediate è quella di specificare il numero degli argomenti opzionali come uno
degli argomenti fissi. Una variazione di questo metodo è l'uso di un argomento
fisso per specificare anche il tipo degli argomenti variabili, come fa la
stringa di formato per \func{printf} (vedi sez.~\ref{sec:file_formatted_io}).

Infine una ulteriore modalità diversa, che può essere applicata solo quando il
tipo degli argomenti lo rende possibile, è quella che prevede di usare un
valore speciale per l'ultimo argomento, come fa ad esempio \func{execl} che
usa un puntatore \val{NULL} per indicare la fine della lista degli argomenti
(vedi sez.~\ref{sec:proc_exec}).

\index{funzioni!\textit{variadic}|)}

\subsection{Il controllo di flusso non locale}
\label{sec:proc_longjmp}

Il controllo del flusso di un programma in genere viene effettuato con le
varie istruzioni del linguaggio C; fra queste la più bistrattata è il
\instruction{goto}, che viene deprecato in favore dei costrutti della
programmazione strutturata, che rendono il codice più leggibile e
mantenibile. Esiste però un caso in cui l'uso di questa istruzione porta
all'implementazione più efficiente e più chiara anche dal punto di vista della
struttura del programma: quello dell'uscita in caso di errore.

\index{salto~non-locale|(} 

Il C però non consente di effettuare un salto ad una etichetta definita in
un'altra funzione, per cui se l'errore avviene in una funzione, e la sua
gestione ordinaria è in un'altra, occorre usare quello che viene chiamato un
\textsl{salto non-locale}.  Il caso classico in cui si ha questa necessità,
citato sia in \cite{APUE} che in \cite{GlibcMan}, è quello di un programma nel
cui corpo principale vengono letti dei dati in ingresso sui quali viene
eseguita, tramite una serie di funzioni di analisi, una scansione dei
contenuti, da cui si ottengono le indicazioni per l'esecuzione di opportune
operazioni.

Dato che l'analisi può risultare molto complessa, ed opportunamente suddivisa
in fasi diverse, la rilevazione di un errore nei dati in ingresso può accadere
all'interno di funzioni profondamente annidate l'una nell'altra. In questo
caso si dovrebbe gestire, per ciascuna fase, tutta la casistica del passaggio
all'indietro di tutti gli errori rilevabili dalle funzioni usate nelle fasi
successive.  Questo comporterebbe una notevole complessità, mentre sarebbe
molto più comodo poter tornare direttamente al ciclo di lettura principale,
scartando l'input come errato.\footnote{a meno che, come precisa
  \cite{GlibcMan}, alla chiusura di ciascuna fase non siano associate
  operazioni di pulizia specifiche (come deallocazioni, chiusure di file,
  ecc.), che non potrebbero essere eseguite con un salto non-locale.}

Tutto ciò può essere realizzato proprio con un salto non-locale; questo di
norma viene realizzato salvando il contesto dello \textit{stack} nel punto in
cui si vuole tornare in caso di errore, e ripristinandolo, in modo da tornare
quando serve nella funzione da cui si era partiti.  La funzione che permette
di salvare il contesto dello \textit{stack} è \funcd{setjmp}, il cui prototipo
è:

\begin{funcproto}{ 
\fhead{setjmp.h}
\fdecl{int setjmp(jmp\_buf env)}
\fdesc{Salva il contesto dello \textit{stack}.} 
}
{La funzione ritorna $0$ quando è chiamata direttamente ed un valore diverso
  da zero quando ritorna da una chiamata di \func{longjmp} che usa il contesto
  salvato in precedenza.}
\end{funcproto}
  
Quando si esegue la funzione il contesto corrente dello \textit{stack} viene
salvato nell'argomento \param{env}, una variabile di tipo
\typed{jmp\_buf}\footnote{anche questo è un classico esempio di variabile di
  \textsl{tipo opaco}.}  che deve essere stata definita in precedenza. In
genere le variabili di tipo \type{jmp\_buf} vengono definite come variabili
globali in modo da poter essere viste in tutte le funzioni del programma.

Quando viene eseguita direttamente la funzione ritorna sempre zero, un valore
diverso da zero viene restituito solo quando il ritorno è dovuto ad una
chiamata di \func{longjmp} in un'altra parte del programma che ripristina lo
\textit{stack} effettuando il salto non-locale. Si tenga conto che il contesto
salvato in \param{env} viene invalidato se la funzione che ha chiamato
\func{setjmp} ritorna, nel qual caso un successivo uso di \func{longjmp} può
comportare conseguenze imprevedibili (e di norma fatali) per il processo.
  
Come accennato per effettuare un salto non-locale ad un punto precedentemente
stabilito con \func{setjmp} si usa la funzione \funcd{longjmp}; il suo
prototipo è:

\begin{funcproto}{ 
\fhead{setjmp.h}
\fdecl{void longjmp(jmp\_buf env, int val)}
\fdesc{Ripristina il contesto dello stack.} 
}
{La funzione non ritorna.}   
\end{funcproto}

La funzione ripristina il contesto dello \textit{stack} salvato da una
chiamata a \func{setjmp} nell'argomento \param{env}. Dopo l'esecuzione della
funzione il programma prosegue nel codice successivo alla chiamata della
\func{setjmp} con cui si era salvato \param{env}, che restituirà il valore
dell'argomento \param{val} invece di zero.  Il valore
dell'argomento \param{val} deve essere sempre diverso da zero, se si è
specificato 0 sarà comunque restituito 1 al suo posto.

In sostanza l'esecuzione di \func{longjmp} è analoga a quella di una
istruzione \instr{return}, solo che invece di ritornare alla riga
successiva della funzione chiamante, il programma in questo caso ritorna alla
posizione della relativa \func{setjmp}. L'altra differenza fondamentale con
\instr{return} è che il ritorno può essere effettuato anche attraverso
diversi livelli di funzioni annidate.

L'implementazione di queste funzioni comporta alcune restrizioni dato che esse
interagiscono direttamente con la gestione dello \textit{stack} ed il
funzionamento del compilatore stesso. In particolare \func{setjmp} è
implementata con una macro, pertanto non si può cercare di ottenerne
l'indirizzo, ed inoltre le chiamate a questa funzione sono sicure solo in uno
dei seguenti casi:
\begin{itemize*}
\item come espressione di controllo in un comando condizionale, di selezione o
  di iterazione (come \instruction{if}, \instruction{switch} o
  \instruction{while});
\item come operando per un operatore di uguaglianza o confronto in una
  espressione di controllo di un comando condizionale, di selezione o di
  iterazione;
\item come operando per l'operatore di negazione (\code{!}) in una espressione
  di controllo di un comando condizionale, di selezione o di iterazione;
\item come espressione a sé stante.
\end{itemize*}

In generale, dato che l'unica differenza fra il risultato di una chiamata
diretta di \func{setjmp} e quello ottenuto nell'uscita con un \func{longjmp} è
costituita dal valore di ritorno della funzione, quest'ultima viene usualmente
chiamata all'interno di un una istruzione \instr{if} che permetta di
distinguere i due casi.

Uno dei punti critici dei salti non-locali è quello del valore delle
variabili, ed in particolare quello delle variabili automatiche della funzione
a cui si ritorna. In generale le variabili globali e statiche mantengono i
valori che avevano al momento della chiamata di \func{longjmp}, ma quelli
delle variabili automatiche (o di quelle dichiarate \dirct{register}) sono in
genere indeterminati.

Quello che succede infatti è che i valori delle variabili che sono tenute in
memoria manterranno il valore avuto al momento della chiamata di
\func{longjmp}, mentre quelli tenuti nei registri del processore (che nella
chiamata ad un'altra funzione vengono salvati nel contesto nello
\textit{stack}) torneranno al valore avuto al momento della chiamata di
\func{setjmp}; per questo quando si vuole avere un comportamento coerente si
può bloccare l'ottimizzazione che porta le variabili nei registri
dichiarandole tutte come \direct{volatile}.\footnote{la direttiva
  \direct{volatile} informa il compilatore che la variabile che è dichiarata
  può essere modificata, durante l'esecuzione del nostro, da altri programmi.
  Per questo motivo occorre dire al compilatore che non deve essere mai
  utilizzata l'ottimizzazione per cui quanto opportuno essa viene mantenuta in
  un registro, poiché in questo modo si perderebbero le eventuali modifiche
  fatte dagli altri programmi (che avvengono solo in una copia posta in
  memoria).}

\index{salto~non-locale|)}


% TODO trattare qui le restartable sequences (vedi
% https://lwn.net/Articles/664645/ e https://lwn.net/Articles/650333/) se e
% quando saranno introdotte

\subsection{La \textit{endianness}}
\label{sec:endianness}

\itindbeg{endianness} 

Un altro dei problemi di programmazione che può dar luogo ad effetti
imprevisti è quello relativo alla cosiddetta \textit{endianness}.  Questa è
una caratteristica generale dell'architettura hardware di un computer che
dipende dal fatto che la rappresentazione di un numero binario può essere
fatta in due modi, chiamati rispettivamente \textit{big endian} e
\textit{little endian}, a seconda di come i singoli bit vengono aggregati per
formare le variabili intere (ed in genere in diretta corrispondenza a come
sono poi in realtà cablati sui bus interni del computer).

\begin{figure}[!htb]
  \centering \includegraphics[height=3cm]{img/endianness}
  \caption{Schema della disposizione dei dati in memoria a seconda della
    \textit{endianness}.}
  \label{fig:sock_endianness}
\end{figure}

Per capire meglio il problema si consideri un intero a 32 bit scritto in una
locazione di memoria posta ad un certo indirizzo. Come illustrato in
fig.~\ref{fig:sock_endianness} i singoli bit possono essere disposti in memoria
in due modi: a partire dal più significativo o a partire dal meno
significativo.  Così nel primo caso si troverà il byte che contiene i bit più
significativi all'indirizzo menzionato e il byte con i bit meno significativi
nell'indirizzo successivo; questo ordinamento è detto \textit{big endian},
dato che si trova per prima la parte più grande. Il caso opposto, in cui si
parte dal bit meno significativo è detto per lo stesso motivo \textit{little
  endian}.

Si può allora verificare quale tipo di \textit{endianness} usa il proprio
computer con un programma elementare che si limita ad assegnare un valore ad
una variabile per poi ristamparne il contenuto leggendolo un byte alla volta.
Il codice di detto programma, \file{endtest.c}, è nei sorgenti allegati,
allora se lo eseguiamo su un normale PC compatibile, che è \textit{little
  endian} otterremo qualcosa del tipo:
\begin{Console}
[piccardi@gont sources]$ \textbf{./endtest}
Using value ABCDEF01
val[0]= 1
val[1]=EF
val[2]=CD
val[3]=AB
\end{Console}
%$
mentre su un vecchio Macintosh con PowerPC, che è \textit{big endian} avremo
qualcosa del tipo:
\begin{Console}
piccardi@anarres:~/gapil/sources$ \textbf{./endtest}
Using value ABCDEF01
val[0]=AB
val[1]=CD
val[2]=EF
val[3]= 1
\end{Console}
%$

L'attenzione alla \textit{endianness} nella programmazione è importante, perché
se si fanno assunzioni relative alla propria architettura non è detto che
queste restino valide su un'altra architettura. Inoltre, come vedremo ad
esempio in sez.~\ref{sec:sock_addr_func}, si possono avere problemi quando ci
si trova a usare valori di un formato con una infrastruttura che ne usa
un altro. 

La \textit{endianness} di un computer dipende essenzialmente dalla architettura
hardware usata; Intel e Digital usano il \textit{little endian}, Motorola,
IBM, Sun (sostanzialmente tutti gli altri) usano il \textit{big endian}. Il
formato dei dati contenuti nelle intestazioni dei protocolli di rete (il
cosiddetto \textit{network order}) è anch'esso \textit{big endian}; altri
esempi di uso di questi due diversi formati sono quello del bus PCI, che è
\textit{little endian}, o quello del bus VME che è \textit{big endian}.

Esistono poi anche dei processori che possono scegliere il tipo di formato
all'avvio e alcuni che, come il PowerPC o l'Intel i860, possono pure passare
da un tipo di ordinamento all'altro con una specifica istruzione. In ogni caso
in Linux l'ordinamento è definito dall'architettura e dopo l'avvio del sistema
in genere resta sempre lo stesso,\footnote{su architettura PowerPC è possibile
  cambiarlo, si veda sez.~\ref{sec:process_prctl}.} anche quando il processore
permetterebbe di eseguire questi cambiamenti.

\begin{figure}[!htbp]
  \footnotesize \centering
  \begin{minipage}[c]{\codesamplewidth}
    \includecodesample{listati/endian.c}
  \end{minipage} 
  \normalsize
  \caption{La funzione \samplefunc{endian}, usata per controllare il tipo di
    architettura della macchina.}
  \label{fig:sock_endian_code}
\end{figure}

Per controllare quale tipo di ordinamento si ha sul proprio computer si è
scritta una piccola funzione di controllo, il cui codice è riportato
fig.~\ref{fig:sock_endian_code}, che restituisce un valore nullo (falso) se
l'architettura è \textit{big endian} ed uno non nullo (vero) se l'architettura
è \textit{little endian}.

Come si vede la funzione è molto semplice, e si limita, una volta assegnato
(\texttt{\small 9}) un valore di test pari a \texttt{0xABCD} ad una variabile
di tipo \ctyp{short} (cioè a 16 bit), a ricostruirne una copia byte a byte.
Per questo prima (\texttt{\small 10}) si definisce il puntatore \var{ptr} per
accedere al contenuto della prima variabile, ed infine calcola (\texttt{\small
  11}) il valore della seconda assumendo che il primo byte sia quello meno
significativo (cioè, per quanto visto in fig.~\ref{fig:sock_endianness}, che sia
\textit{little endian}). Infine la funzione restituisce (\texttt{\small 12})
il valore del confronto delle due variabili. 

In generale non ci si deve preoccupare della \textit{endianness} all'interno
di un programma fintanto che questo non deve generare o manipolare dei dati
che sono scambiati con altre macchine, ad esempio via rete o tramite dei file
binari. Nel primo caso la scelta è già stata fatta nella standardizzazione dei
protocolli, che hanno adottato il \textit{big endian} (che viene detto anche
per questo \textit{network order}); vedremo in sez.~\ref{sec:sock_func_ord} le
funzioni di conversione che devono essere usate.

Nel secondo caso occorre sapere quale \textit{endianness} è stata usata nei
dati memorizzati sul file e tenerne conto nella rilettura e nella
manipolazione e relativa modifica (e salvataggio). La gran parte dei formati
binari standardizzati specificano quale \textit{endianness} viene utilizzata e
basterà identificare qual'è, se se ne deve definire uno per i propri scopi
basterà scegliere una volta per tutte quale usare e attenersi alla scelta.

\itindend{endianness}


% LocalWords:  like exec kernel thread main ld linux static linker char envp Gb
% LocalWords:  sez POSIX exit system call cap abort shell diff errno stdlib int
% LocalWords:  SUCCESS FAILURE void atexit stream fclose unistd descriptor init
% LocalWords:  SIGCHLD wait function glibc SunOS arg argp execve fig high kb Mb
% LocalWords:  memory alpha swap table printf Unit MMU paging fault SIGSEGV BSS
% LocalWords:  multitasking text segment NULL Block Started Symbol fill black
% LocalWords:  heap stack calling convention size malloc calloc realloc nmemb
% LocalWords:  ENOMEM ptr uClib cfree error leak smartpointers hook Dmalloc brk
% LocalWords:  Gray Watson Electric Fence Bruce Perens sbrk longjmp SUSv BSD ap
% LocalWords:  ptrdiff increment locking lock copy write capabilities IPC mlock
% LocalWords:  capability MEMLOCK limits getpagesize RLIMIT munlock sys const
% LocalWords:  addr len EINVAL EPERM mlockall munlockall flags l'OR CURRENT IFS
% LocalWords:  argc argv parsing questofile txt getopt optstring switch optarg
% LocalWords:  optind opterr optopt POSIXLY CORRECT long options NdA group
% LocalWords:  option parameter list environ PATH HOME XPG tab LOGNAME LANG PWD
% LocalWords:  TERM PAGER TMPDIR getenv name SVr setenv unsetenv putenv opz gcc
% LocalWords:  clearenv libc value overwrite string reference result argument
% LocalWords:  socket variadic ellipsis header stdarg execl self promoting last
% LocalWords:  float double short register type dest src extern setjmp jmp buf
% LocalWords:  env return if while Di page cdecl  rectangle node anchor west PS
% LocalWords:  environment rounded corners dashed south width height draw east
% LocalWords:  exithandler handler violation inline SOURCE SVID XOPEN mincore
% LocalWords:  length unsigned vec EFAULT EAGAIN dell'I memalign valloc posix
% LocalWords:  boundary memptr alignment sizeof overrun mcheck abortfn enum big
% LocalWords:  mprobe DISABLED HEAD TAIL touch right emacs OSTYPE endianness IBM
% LocalWords:  endian little endtest Macintosh PowerPC Intel Digital Motorola
% LocalWords:  Sun order VME  loader Windows DLL shared objects PRELOAD termios
% LocalWords:  is to LC SIG str mem wcs assert ctype dirent fcntl signal stdio
% LocalWords:  times library utmp syscall number Filesystem Hierarchy pathname
% LocalWords:  context assembler sysconf fork Dinamic huge segmentation program
% LocalWords:  break store using intptr ssize overflow ONFAULT faulting alloc
%  LocalWords:  scheduler pvalloc aligned ISOC ABCDEF

%%% Local Variables: 
%%% mode: latex
%%% TeX-master: "gapil"
%%% End: 
